% concatenated from main.tex
\documentclass[11pt,reqno,a4paper]{amsart}
\usepackage{amsmath,amsfonts,amssymb,amsthm,amsxtra,latexsym,amscd,enumerate,verbatim}
\usepackage[margin=2.0cm]{geometry}
\usepackage[abbrev,lite,nobysame]{amsrefs}
\usepackage{graphics,graphicx}
\usepackage{subcaption}
\usepackage{tabularx}
\usepackage[usenames,dvipsnames]{color}
\usepackage{times}
\usepackage{bbm}
\usepackage[colorlinks=true, pdfstartview=FitV, linkcolor=blue, citecolor=blue, urlcolor=blue]{hyperref}
\usepackage{tikz}
\usepackage{enumerate,enumitem}
\usepackage{mathtools}
\makeatletter
\usepackage[utf8]{inputenc}
\usepackage{amsmath,amsfonts,amssymb,amsxtra,latexsym,amscd,enumerate,amsthm,verbatim}
\usepackage[abbrev,lite,nobysame]{amsrefs}
\usepackage{float}
\usepackage{tikz}
\usepackage{pgfplots}
\usepackage{caption}
\usepackage[colorlinks=true]{hyperref}
\hypersetup{urlcolor=blue, citecolor=blue}
\usepackage{hyperref}
\numberwithin{equation}{section}
\def\dd{{\mathrm{d}}}
\newcommand{\bl}{\textcolor{blue}}
\newcommand{\R}{\mathbb{R}}
\newcommand{\N}{\mathbb{N}}
\newcommand{\C}{\mathbb{C}}
\newcommand{\Z}{\mathbb{Z}}
\newcommand{\defeq}{\mathrel{\mathop:}=}
\newcommand{\eqdef}{\mathrel{\mathop=}:}
\newcommand{\tg}{{\tt g}}
\newcommand{\inn}{{\quad\hbox{in } }}
\newcommand{\ve}{\varepsilon}
\newcommand{\be}{\begin{equation}}
\newcommand{\ee}{\end{equation}}
\renewcommand{\div}{\mathop{\rm div}}
\newtheorem{definition}{Definition}
\newtheorem{lemma}{Lemma}[section]
\newtheorem{prop}{Proposition}[section]
\newtheorem{theorem}{Theorem}
\newtheorem{corollary}{Corollary}[section]
\newtheorem{remark}{Remark}[section]
\newtheorem{claim}{Claim}
\pgfplotsset{compat=1.18} 

\title[Clustered vortex helices with compactly supported cross-sectional vorticity]{Clustered vortex helices with compactly supported cross-sectional vorticity in the 3D Euler equations}
\author[A.~Averkiou]{Averkios Averkiou}
\address{\noindent A.~Averkiou: Department of Mathematical Sciences University of Bath, Bath BA2 7AY,
United Kingdom.}
\email{aa4119@bath.ac.uk}

\author[M.~Musso]{Monica Musso}
\address{\noindent M.~Musso: Department of Mathematical Sciences University of Bath, Bath BA2 7AY,
United Kingdom.}
\email{m.musso@bath.ac.uk}
\author[F.~Yu]{Fang Yu}
\address{\noindent F.~Yu: School of Mathematical Sciences, Chongqing University of Technology,
Chongqing, P.R. China}
\email{yufang@cqut.edu.cn}

\begin{document}
\begin{abstract}
We consider the three-dimensional incompressible Euler equations for helical flows without swirl. By adapting gluing techniques, we construct the first smooth multi-vortex solution in the whole space $\mathbb{R}^3$ exhibiting a cluster of collapsing helical filaments, with the associated cross-sectional vorticity remaining compactly supported in $\mathbb{R}^2$
 for all times. Our result generalises previous collapsing configurations in $\mathbb{R}^3$ with rapidly decaying vorticity cores, and extends related variational solutions obtained in infinite cylindrical domains.
\end{abstract}
\maketitle
\section{Introduction}
The motion of a three-dimensional inviscid and incompressible fluid without external forcing is described by the Euler equations in $\R^3$. For a sufficiently smooth velocity field $\mathbf{u}: \R^3\times [0,T) \to \R^3$ with $\nabla\cdot \mathbf{u} =0,$ vorticity $\vec{\omega}\defeq \nabla\times{\mathbf{u}}: \R^3\times [0,T) \to \R^3,$ stream function $\vec{\psi}: \R^3 \times [0,T) \to \R^3$ and prescribed initial data $\mathbf{u}(\cdot,0)=u_0(x)$, the Euler system in vorticity formulation reads
\begin{equation}
\label{euler}
\begin{aligned}
\partial_{t}\vec \omega  +
(\mathbf{u}\cdot \nabla){\vec \omega}
&=( \vec \omega \cdot \nabla)\mathbf {u}  &&  \inn \, \R^3\times (0,T), \\
\quad \mathbf{u}  = \nabla \times \vec {\psi},\ &
-\Delta \vec \psi =  \vec \omega  && \inn \, \R^3\times (0,T),\\
{\vec \omega}(\cdot,0)  &=  \nabla \times u_0 && \inn \, \R^3.
\end{aligned}
\end{equation}
We are interested in the rigorous construction of smooth solutions to \eqref{euler} with vortex filament structure, where the vorticity remains highly concentrated around a smooth curve evolving over time. Despite the growing interest in the vortex dynamics of such flows in recent years, their study dates back to the foundational works of Helmholtz and Kelvin \cites{MR1579057,thomson1868vi}. Formal asymptotic considerations suggest that when the vorticity is concentrated in a thin tubular neighbourhood of radius $\ve>0$ around a smooth curve $\mathcal{G}(t),$ the associated filament evolves to leading order according to the binormal flow with velocity  $\mathcal{O}\left(|\log\ve|\right)$. More precisely, for an arclength parametrisation $\gamma(s,t)$ and a circulation constant $\bar c$, the formal derivations in \cites{da1906sul, levi1908sull}  indicate that the evolution is governed by
\[\partial_{t}\gamma=\bar{c}|\log\ve|\left(\gamma_{s}\times\gamma_{ss}\right)\ \quad \text{as}\quad \ve\to 0,\]
which one can also express as
\begin{equation}\label{binormal}
\partial_{\tau}\gamma=\bar{c}\varkappa\mathbf{B}_{{\mathcal{G}(\tau)}}, \quad t = |\log\ve|^{-1}\tau,
\end{equation}
with $\varkappa$ denoting the curvature and $\mathbf{B}_{\mathcal{G}(\tau)}$ the binormal unit vector. In this regime, the \emph{vortex filament conjecture} concerns the rigorous verification that concentrated vorticity structures remain localised and evolve at leading order along the trajectories determined by \eqref{binormal}. Beyond the conditional result in \cite{MR3609248}, significant progress has been achieved in the case of circles and helices, which constitute symmetric solutions of \eqref{binormal}.

For circular vortex filaments, the validity of the binormal law \eqref{binormal} was first established in \cite{fraenkel1970steady} through the construction of \emph{vortex rings} propagating with constant speed $\mathcal{O}\left(|\log\ve|\right)$, with subsequent  contributions including \cites{MR3101789,MR422916,MR1738349,MR4404610, guo2026global}. In the helical framework, the vortex filament conjecture was rigorously verified in \cite{MR4417384}, where the authors constructed \emph{vortex helices}; see also \cites{qin2024concentrated,MR4555251,MR4534710,MR5036920} and the references therein. Of particular relevance to the present work are the helical multi-vortex solutions in \cites{MR4899797,cao2023clustered,MR5030554,cao2025co,donati2024dynamics}, covering among others nearly parallel and collapsing configurations, while results on global well-posedness and long-time dynamics appear in  \cites{MR4809294,guo2024long,MR3320529}. Finally, helical structures have also been investigated in related models, such as the Ginzburg–Landau  \cites{MR4493621,DUAN2026114325}, the Gross–Pitaevskii  \cites{MR2181460,MR4540752}, and the Schrödinger map \cite{MR3449250} equations. 
\subsection{Statement of the main result}
This paper is concerned with the rigorous construction of the first smooth global-in-time solution $\vec{\omega}_{\ve}(x,\tau)$ to \eqref{euler} in the whole space $\R^3,$ which concentrates, as $\ve \to 0,$ near multiple collapsing helices undergoing coupled rotation and translation at nearly identical speeds, with the main novelty that the associated cross-sectional vorticity remains compactly supported in $\R^2$ for all times. We refer to such a solution as a \emph{cluster of vortex helices with compactly supported cross-sectional vorticity}.

In this regard, given an integer $N\geq 1$ and $h>0,$ we take planar points $P_i\defeq(a_i,b_i),\,i=1,\dots,N,$ and consider the evolving circular helices  $\mathcal{G}_i(\tau)$ of radius $R_i=\sqrt{a_i^2+b_i^2}$ and common pitch $h>0$, parametrised by arclength as 
\begin{equation} \label{parhelix}
	\gamma_i(s,\tau ) = \left( \begin{matrix}
		a_i \cos \left(
		\frac{s-\sigma_{1,i} \tau}{\sqrt{h^2 + R_i^2}}\right)-b_i\sin \left(\frac{s-\sigma_{1,i}\tau}{\sqrt{h^2+R_i^2}} \right)
		\\
		a_i \sin \left(
		\frac{s-\sigma_{1,i} \tau}{\sqrt{h^2 + R_i^2}}\right)+b_i\cos \left(
		\frac{s-\sigma_{1,i} \tau}{\sqrt{h^2 + R_i^2}}\right)
		\\
		\frac{ h s + \sigma_{2,i} \tau }{\sqrt{h^2+ R_i^2}}
	\end{matrix} \right) \in \R^3, \quad \sigma_{1,i} = \frac{\bar{c}_i\,  h}{R_i^2 + h^2},
	\quad 
	\sigma_{2,i} = \frac{ \bar{c}_i\, R_i^2}{R_i^2 + h^2}.
	\end{equation}
One can readily verify that each $\gamma_i$ has curvature $\frac{R_i}{R_i^2+h^2}$, torsion $\frac{h}{R_i^2 + h^2}$, and evolves according to the binormal flow \eqref{binormal} with $\bar{c}_i=\bar{c}\kappa_i$, where $\kappa_i$ is the circulation and $\bar{c}$ is determined by the regularisation profile; see \eqref{defGamma}. Hereafter, we identify the base point  $(a_i,b_i,0)$ of each helix with the planar point $(a_i,b_i)$.

In view of our interest in  colliding helical filaments, for a fixed $r_0>0$ we consider points $P_i=(a_i,b_i)$ in \eqref{parhelix}  of the specific form 
\begin{equation}\label{formpoints}
P_i=(r_0+\tilde{r},0)+ \frac{\widehat{P}_i}{|\log\ve|}, \quad |\tilde{r}|\leq \delta \frac{\log|\log\ve|}{|\log\ve|}, \quad\delta<|\widehat{P}_i|<\delta^{-1}, \quad i=1,\dots,N,
\end{equation}
where $\delta>0$ is sufficiently small and independent of $\ve>0.$ Since $P_i=(a_i,b_i)\to (r_0,0)$ as $\ve\to 0$ for all $i\in\{1,\dots,N\},$ we infer that the helices $\mathcal{G}_i(\tau)$ in \eqref{parhelix} collapse onto the same limiting helix as $\ve \to 0.$ 

Furthermore, we assume that the points $\widehat{P}_i$ in \eqref{formpoints} are uniformly separated, namely there exists a constant $d>0$ independent of $\ve>0$ such that 
\begin{equation}\label{unifdistance}
\min_{i\neq j} |\widehat{P}_j-\widehat{P}_i|\geq d,
\end{equation}
and we also impose a reflection symmetry on the set $\{\widehat{P}_1,\dots,\widehat{P}_N\}$ with respect to the first coordinate axis, i.e.
\begin{equation}\label{symmetryforpoints}
\mathcal{P}=(\mathcal{P}_{1},\mathcal{P}_{2}) \in \{\widehat{P}_1,\dots,\widehat{P}_N\} \iff (\mathcal{P}_{1},-\mathcal{P}_{2}) \in \{\widehat{P}_1,\dots,\widehat{P}_N\}.\end{equation} 
Finally, for $h>0$ and  $\kappa_1,\dots,\kappa_N\in\R$ such that $\sum\limits_{i=1}^{N}\kappa_i>0$, the points $\widehat{P}_i$ admit the decomposition
\begin{equation}\label{points1}
\widehat{P}_i=P_i^{b}+q_i, \quad i=1,\dots,N,
\end{equation}
so that if we write in anisotropic coordinates 
\begin{equation}\label{rescaledmainpoints}
{P}^{b}_i=\left(\frac{h\,\widetilde{P}^{b}_{i,1}}{\sqrt{h^2+r_0^2}},\widetilde{P}^{b}_{i,2}\right), \quad \widetilde{P}^{b}_i=\left(\widetilde{P}^{b}_{i,1},\widetilde{P}^{b}_{i,2}\right),
\end{equation}
the configuration $\left(\widetilde{P}^{b}_{1},\dots,\widetilde{P}^{b}_{N}\right)$ solves, for all $i=1,\dots,N,$ the system of balancing conditions
\begin{equation}\label{balancingcond}
\begin{aligned}
\sum_{j \neq i} \kappa_j\frac{\widetilde{P}^{b}_{i,1}-\widetilde{P}^{b}_{j,1}}{|\widetilde{P}^{b}_i-\widetilde{P}^{b}_j|^2}&=\left(\frac{\kappa_ih r_0}{2(h^2+r_0^2)^{\frac 3 2}}+\alpha\frac{h r_0}{\nu'(1)\sqrt{h^2+r_0^2}}\right),  \\
\sum_{j \neq i} \kappa_j\frac{\widetilde{P}^{b}_{i,2}-\widetilde{P}^{b}_{j,2}}{|\widetilde{P}^{b}_i-\widetilde{P}^{b}_j|^2}&=0,
\end{aligned}
\end{equation}
for some constants $\alpha>0$ and $\nu'(1)<0$; see \eqref{defGamma}. Here and throughout this work, the system \eqref{balancingcond} uniquely determines $\alpha>0$ as  
\begin{equation}\label{uniformspeed}
\alpha=\frac{-\nu'(1)}{2(h^2+r_0^2)}\frac{\sum\limits_{i=1}^{N} \kappa_i^2}{\sum\limits_{i=1}^{N}\kappa_{i}}.
\end{equation}
In addition, the terms ${q}_i$ in \eqref{points1} are lower order perturbations, with $|q_i|\to 0$ as $\ve \to 0$. 

Notably, admissible configurations of points $P^{b}=\left(P^{b}_1,\dots,P^{b}_N\right)$ satisfying \eqref{unifdistance}, \eqref{symmetryforpoints}, \eqref{rescaledmainpoints} and \eqref{balancingcond} are exhibited in \cites{MR4404610,MR4350893}. In this context, since \eqref{balancingcond} is translation-invariant, we call a solution $P^{b}$ as in \eqref{rescaledmainpoints} \emph{nondegenerate} if the linearisation of \eqref{balancingcond} at $P^{b}$ has a one-dimensional kernel, arising from the symmetry assumption  \eqref{symmetryforpoints}. A more detailed analysis of this condition and its consequences is deferred to Section \ref{reduced}, where it plays a crucial role in locating the points $\widehat{P}_i$ in \eqref{points1} as small perturbations of $P_i^{b}$,\, $i=1,\dots,N,$ thereby determining the adjusted vortex centers of the helices associated with an exact smooth solution of \eqref{euler}.

The following Theorem is the main result of this work.
\begin{theorem}\label{maintheorem}
Fix an integer $N\geq 1$, constants $r_0>0,\,h>0,$ and  $\kappa_1,\dots,\kappa_N\in\R$ satisfying $\sum\limits_{i=1}^{N}\kappa_i>0.$ For each $i=1,\dots,N$, denote by $\mathcal{G}_i(\tau)$ the evolving helix parametrised by \eqref{parhelix}, centered at the point $P_i=(a_i,b_i)$ in \eqref{formpoints}, and let $\delta_{\mathcal{G}_i(\tau)}$ designate a uniform Dirac measure supported on $\mathcal{G}_i(\tau)$ and $\mathbf{T}_{\mathcal{G}_i(\tau)}$ its tangent unit vector. Moreover, assume  
that the configuration $P^{b}=\left(P_1^{b},\dots,P_N^{b}\right)$ satisfying \eqref{unifdistance},  \eqref{symmetryforpoints} and \eqref{rescaledmainpoints}, is a nondegenerate solution of  \eqref{balancingcond}. Then, there exist 
$\tilde{r}_{*}\in\R$,  $q_1,\dots,q_N\in\R^2$ and specific points of the form
\begin{equation}\label{pointstheorem}
P_i=(r_0+\tilde{r}_{*},0)+\frac{1}{|\log\ve|}\widehat{P}_i,\quad \widehat{P}_i=P_{i}^{b}+q_i, \quad |\tilde{r}_{*}|, |q_i| \lesssim \frac{\log|\log\ve|}{|\log\ve|}, \quad i=1,\dots,N,
\end{equation}
which determine a smooth global-in-time solution $\vec{\omega}_{\ve}(x,\tau)$ of \eqref{euler} with cross-sectional vorticity compactly supported in $\R^2$ for all times, such that for all $\tau\in\R$ it satisfies
\[\vec{\omega}_{\ve}(x,|\log\ve|^{-1}\tau) \rightharpoonup \sum_{i=1}^{N}\kappa_i\delta_{\mathcal{G}_i(\tau)}\mathbf{T}_{\mathcal{G}_i(\tau)}\quad \mbox{as} \quad \ve\to 0,\]
in the sense of measures. 
\end{theorem}
We now present some remarks on Theorem \ref{maintheorem} and its relation to the existing literature.
\begin{remark}
Theorem \ref{maintheorem} provides the first smooth multi-vortex helical solution $\vec{\omega}_{\ve}(x,\tau)$ of \eqref{euler} in the whole space $\R^3$ whose vorticity concentrates near collapsing filaments as $\ve \to 0$, while the associated cross-sectional vorticity remains compactly supported in $\R^2$ for all times. More precisely, for all  $\ve>0$ sufficiently small, $\alpha>0$ given by \eqref{uniformspeed} and $P_i$ as in \eqref{pointstheorem}, it holds
\begin{equation}\label{remark1prop}
\operatorname{supp}\vec{\omega}_{\ve}(\cdot,x_3,\tau)\subset \bigcup_{i=1}^{N} B_{\ve|\log\ve|^2}\left(R_{\frac{x_3}{h}-\alpha\tau}P_i\right),
\quad R_{\theta}=\begin{pmatrix}\cos\theta&-\sin\theta\\\sin\theta&\cos\theta\end{pmatrix},
\end{equation}
for all $(x_3,\tau)\in\R\times\R,$ hence each planar cross-section of $\vec{\omega}_{\ve}(x,\tau)$ is supported in a union of $N$ disjoint compact balls. Consequently, for every $\tau\in\R$, the solution $\vec{\omega}_{\ve}(x,\tau)$ remains confined within $N$ non-overlapping helical tubes of radius $\ve|\log\ve|^2$, whose vortex cores collapse onto a single helix as $\ve \to 0$.
\end{remark}
\begin{remark}
A related multi-vortex helical solution of \eqref{euler} in bounded domains is obtained in \cite{donati2024dynamics}, where the authors consider vorticity initially concentrated near helices of distinct radii and show that the concentration persists in a finite time interval $\tau \in [0,T]$, with cross-sectional vorticity supported in a thin annulus. Nevertheless, their filaments remain uniformly separated at distance $\mathcal{O}(1)$, thus the interaction between filaments is asymptotically negligible compared to self-induced motion, effectively reducing the problem to a single-helix setting. In contrast, our filament separation scale is $\mathcal{O}(|\log\ve|^{-1})$, placing the problem in a strongly coupled interaction regime, while concentration persists for all $\tau \in\R$. For a similar reason related to filament separation, the polygonal solution of \cite{MR5036920} does not fit the current framework. Moreover, \cite{cao2023clustered} is concerned with clustered helical solutions in infinite cylinders obtained through variational methods. However, their analysis relies crucially on a $C^{1}$ expansion of the Green’s function of a uniformly elliptic operator in bounded planar domains. Since this expansion fails in the whole $\R^2$ due to the loss of uniform ellipticity, their approach does not seem to be directly adaptable to our setting.
\end{remark}

\subsection{Helical symmetry with no swirl reduction} 
Our construction is carried out within the class of helically symmetric flows with no helical swirl \cites{MR2505860,MR1717127}, where the Euler dynamics in \eqref{euler} reduces to a transport-elliptic system in $\R^2$. More precisely, writing $x=(x',x_3)\in\R^3$, $x'=(x_1,x_2)\in\R^2$ and  $t=|\log\ve|^{-1}\tau$, we consider the pair $\left(w(x',\tau),\psi(x',\tau)\right)$ solving
\begin{equation}\label{PB}
		\begin{aligned}
			\left\{
			\begin{aligned}
				|\log\ve|\partial_{\tau}\omega + \nabla^\perp \psi \cdot \nabla w &=0 && {\mbox {in}} \quad \R^2 \times (-\infty, +\infty),
				\\
				- \nabla\cdot (K \nabla \psi) &= w && {\mbox {in}} \quad \R^2 \times (-\infty,+\infty),
			\end{aligned}
			\right.
		\end{aligned}
	\end{equation}
	with $(a,b)^\perp = (b,-a)$ and $K (x_1,x_2) $ denoting the coefficient matrix
\[
	K(x_1 , x_2 ) = \frac{1}{h^2+x_1^2+x_2^2}
	\left(
	\begin{matrix}
		h^2+x_2^2 & -x_1  x_2\\
		-x_1 x_2 & h^2+x_1^2
	\end{matrix}
	\right).\]
If $w(x',\tau)$ solves \eqref{PB}, then the vorticity vector
	\begin{align}
\label{omegaHelical}
		\vec \omega(x,\tau)  =  \frac{1}{h} w\left( R_{-\frac {x_3}h} x',\tau\right)  \, \begin{pmatrix}-x_2\\x_1\\h\end{pmatrix}, \quad R_{\theta}=\begin{pmatrix} \cos\theta& -\sin\theta\\ \sin\theta& \cos\theta \end{pmatrix},
	\end{align}
	is helically symmetric with no swirl and satisfies \eqref{euler}. Moreover, seeking rotating solutions of the form
\begin{equation}
\label{rotansatz}
w (x', \tau) = W \left( R_{ \alpha \tau } x' \right), \quad \psi (x',\tau) = \Psi \left( R_{ \alpha \tau } x' \right), \quad x'=(x_1,x_2)\in\R^2,
\end{equation}
reduces \eqref{PB} to the semilinear elliptic equation
\begin{equation}\label{semilin}
\nabla_{\tilde{x}}\cdot\left(K(\tilde x)\nabla_{\tilde{x}}\Psi\right)+F\left(\Psi-\frac{\alpha}{2}|\log\ve||\tilde{x}|^2\right)=0 \quad \mbox{in} \quad \R^2, \quad \tilde{x}=R_{\alpha\tau}x'.
\end{equation}
Our aim is to find a suitable smooth nonlinear function $F$ and a stream function $\Psi$ for \eqref{semilin}, so that the 
vorticity $W(\tilde x)\defeq F\left(\Psi-\frac{\alpha}
{2}|\log\ve||\tilde x|^2\right)$
satisfies
\[W(\tilde x)\sim \sum_{i=1}^{N}\kappa_i\delta_{P_i},\quad \operatorname{supp}W(\tilde x) \subset \bigcup_{i=1}^{N} B_{\tilde{\rho}_{\ve}}(P_i) , \]
with $P_i$ as in \eqref{formpoints} and $\tilde{\rho}_{\ve}=o\left(|\log\ve|^{-1}\right)$ as $\ve\to 0$. In this way, the cross-sections are pairwise disjoint since $\min\limits_{i\neq j}|P_i-P_j|\geq\frac{d}{|\log\ve|}$, and the profile  $w(x',\tau)=W(R_{\alpha\tau} x')$ in \eqref{rotansatz}  
is rigidly rotating with angular speed $\alpha$ as in \eqref{uniformspeed}, while remaining compactly supported for all times in the rotated union of planar vortex cores. Consequently, due to \eqref{omegaHelical} we obtain a smooth $3$D multi-filament solution $\vec{\omega}(x,\tau)$, which is confined in thin non-overlapping helical tubes globally in time.
\begin{remark}\label{decoupling}
Theorem \ref{maintheorem} is proved via an Inner–Outer gluing scheme, in the spirit of \cite{MR4899797}. A key novelty of the present work is the choice of nonlinearity $F$ in \eqref{semilin}, motivated by the equation $\Delta u + u^{\gamma}_{+}=0 \,$ in $\,\R^2$, where $u_{+}\defeq\max(0,u)$ and $\gamma>3$. This setting differs signifincatly from that in \cite{MR4899797}, where $F(s)\sim e^{s}$. While the latter yields rapidly decaying cross-sectional vorticity, the current choice $F(s)\sim s_{+}^{\gamma}$ gives rise to a compactly supported one, as in \eqref{remark1prop}. Moreover, since our construction involves the operator  $\Delta+\gamma\Gamma_{+}^{\gamma-1}$ in $\R^2$, with $\Gamma$ given in \eqref{defGamma}, the presence of the radial kernel element $Z_{0}=\mathcal{O}(\log(2+|y|))$ in \eqref{elemkernel} poses a major technical difficulty, since the quantity $
\int_{\R^{2}} \gamma \Gamma_{+}^{\gamma-1} Z_{0} \neq 0$ induces a strong coupling in the Inner-Outer system (see Proposition \ref{propsec8}). In contrast, the Liouville operator $\Delta+e^{u}$ employed in \cite{MR4899797}, with $\Delta u + e^{u}=0$ in $\R^2$, has a radial kernel element $\widetilde {Z}_{0}=\mathcal{O}(1)$ with $\int_{\R^{2}} e^{u}\widetilde {Z}_{0}=0$, hence the corresponding Inner–Outer system is sufficiently decoupled. To resolve this, we refine the gluing procedure in two main directions. First, we introduce perturbations $\mu_j^{*}$ in the scaling parameters of Proposition \ref{propscaling} to weaken the coupling and obtain solvability of the projected problem in Proposition \ref{propsec8}. Second, to make this modification compatible with the overall gluing scheme, we adapt the solution procedure of the Outer problem (see section \ref{subsecouter}) so that these scaling perturbations are a posteriori sufficiently small in $\ve>0.$
\end{remark}
Following the preceding discussion, in the sequel we write $x$ instead of $\tilde{x}$ in \eqref{semilin} for ease of notation, and dedicate the rest of the paper to solving  the problem
\begin{equation}\label{eqtosolve}
\nabla_{x}\cdot\left(K\nabla_{x}\Psi\right)+F\left(\Psi-\frac{\alpha}{2}|\log\ve||x|^2\right)=0 \quad \mbox{in} \quad \R^2.
\end{equation}
As explained earlier, to establish Theorem \ref{maintheorem} it suffices to select a suitable nonlinearity $F$ and construct a stream function  $\Psi$ for \eqref{eqtosolve}, so that the induced vorticity  $W(x)\defeq F\left(\Psi-\frac{\alpha}{2}|\log\ve||x|^2\right)$ has the properties
\begin{equation}\label{propvorticity}
 F\left(\Psi-\frac{\alpha}{2}|\log\ve||x|^2\right)\sim \sum_{i=1}^{N} \kappa_i \delta_{P_i}, \quad \operatorname{supp}W(x) \subset \bigcup_{i=1}^{N} B_{\ve|\log\ve|^2}(P_i),
 \end{equation}
 where $P_i=(a_i,b_i)\to (r_0,0)$ as $\ve\to0\,$ for all $i=1,\dots,N,$ as can be verified by \eqref{formpoints}. 
\subsection{Structure of the paper}
The organisation of this work is as follows. In Section \ref{Section 2}, we construct a local regularised approximation of the Green’s function of $\nabla_{x}\cdot(K\nabla_{x}\cdot)$ in $\R^2$, based on a radial profile $\Gamma$ solving a semilinear elliptic equation and on a suitable change of variables. Section \ref{section3} extends this construction to a global multi-vortex approximation, while Section \ref{Section 4} is devoted to the properties of the associated scaling parameters. In Section \ref{Section 5}, we choose a compactly supported nonlinearity $F$ and derive the corresponding error estimates. Section \ref{Section 6} introduces the \emph{Inner-Outer gluing scheme} and describes the strategy for perturbing the approximation into an exact solution, relying on the linear theories of Section \ref{Section 7}. Finally, in Section \ref{Section projected} we solve a projected linearised problem, and in Section \ref{reduced} we suitably adjust the vortex centers to obtain a genuine solution.
\section{Local approximate stream function}\label{Section 2} 
Let $h>0 $ and consider a point $P=(a,b)\in\R^2\setminus\{0\}$. This section is concerned with the construction of an approximate  stream function $\Psi$ with the property 
\begin{equation}\label{ori}
-\nabla_{x}\cdot (K\nabla_{x}\Psi) \sim c\delta_P
\end{equation}
locally around $P$, with
\begin{equation}\label{operL}
L_x=\nabla_{x}\cdot (K\nabla_{x}\cdot)
\quad \text{and}\ \quad
K=\frac{1}{h^2+x_1^2+x_2^2}
\begin{pmatrix}
h^2+x_2^2 & -x_1x_2 \\
-x_1x_2 & h^2+x_1^2
\end{pmatrix}.
\end{equation}
As a preliminary step, we analyse the behaviour of $L_x$ in a small neighbourhood of the point $P$, which will enable us to derive a local approximation of the Green’s function for \eqref{operL}. 

Following \cite{MR4899797}, we introduce the linear change of variables
\begin{equation}\label{var}
x-P=A[P]z,\quad
A[P]=
\begin{pmatrix}
\frac{ah}{R\sqrt{h^2+R^2}}
& -\frac{b}{R}\\
\frac{bh}{R\sqrt{h^2+R^2}}
& \frac{a}{R}
\end{pmatrix},
\end{equation}
which normalises the coefficient matrix $K$ in \eqref{operL} at the point $P=(a,b)$.

The following Proposition holds.
\begin{prop}\label{expansionofL}
Let $z$ be the variable defined in \eqref{var}. Then, the operator $L_x$ in \eqref{operL} can be expressed as  
\[L_{x}=\Delta_{z}+B_0,\]
where
\begin{align*}
B_{0}
&= \left(\frac{h^2(R^2-r^2)+z_2^2(h^2+R^2)}{h^2(h^2+r^2)}\right)\partial_{z_1z_1}
+\frac1{(h^2+r^2)}\left(\left(z_1\frac h{\sqrt{h^2+R^2}}+R\right)^2-r^2\right)\partial_{z_2z_2} \nonumber\\
 &\hspace{5mm}- 2\frac{\sqrt{h^2+R^2}}{h(h^2+r^2)}z_2\left(z_1\frac{h}{\sqrt{h^2+R^2}}+R\right)
\partial_{z_1z_2}
 \nonumber\\
&\hspace{5mm}-\frac{z_1(h^2+R^2)+Rh\sqrt{h^2+R^2}}{h^2(h^2+r^2)}\left(1+\frac{2h^2}{h^2+r^2}\right)
\partial_{z_1}
 - \frac{z_2}{h^2+r^2}\left(1+\frac{2h^2}{h^2+r^2}\right)\partial_{z_2},
\end{align*}
and
\begin{equation}\nonumber
r^2=|x|^2=R^2+2R\frac h{\sqrt{h^2+R^2}}z_1 + \frac{h^2}{h^2+R^2}z_1^2+z_2^2.
\end{equation}
\end{prop}

\begin{proof}
See \cite{MR4899797}*{Proposition 2.1}.
\end{proof}

Given any small $\delta>0$, Proposition \ref{expansionofL} shows that in the bounded region $|z|<\delta$, the operator $L_x$ in \eqref{operL} may be regarded as the Laplace operator $\Delta_z$ up to lower-order perturbations. 

More precisely, we have
\begin{equation}\label{oper}
L_x=\Delta_z+B_{0},
\end{equation}
where $B_{0}$ admits the expansion 
\begin{equation}
\begin{aligned}\nonumber
B_0
&= \left(-2\frac{Rh}{(h^2+R^2)^{\frac 3 2}}z_1+O(|z|^2)\right)\partial_{z_1z_1}
+\mathcal{O}(|z|^2)\partial_{z_2z_2} -\left(2\frac {R}{h\sqrt{h^2+R^2}}z_2+\mathcal{O}(|z|^2)\right)\partial_{z_1z_2}  \nonumber \\
&-\left(\frac{R}{h\sqrt{h^2+R^2}}\left(1+\frac{2h^2}{h^2+R^2}\right)+\mathcal{O}(|z|)\right)
\partial_{z_1} -\left(\frac{z_2}{h^2+R^2}\Big(1+\frac{2h^2}{h^2+R^2}\right)
+\mathcal{O}\left(|z|^2)\right)
\partial_{z_2}.
\end{aligned}
\end{equation}
Furthermore, \eqref{ori} is now formulated as
\begin{equation}\label{oriz}
-(\Delta_z+B_0)\psi = c\,\delta_{0}, 
\qquad 
\psi(z)=\Psi\bigl(P+A[P]z\bigr),
\end{equation}
where $\delta_{0}$ represents a Dirac delta at the origin and $A[P]$ is the matrix employed in the change of variables \eqref{var}. Consequently, an approximate regularisation of the Green’s function for $L_x$ can be obtained through a regularisation of the Green's function of $\Delta_z$. 

Motivated by this observation, for $\gamma>3$ and $s_{+}\defeq\max\{0,s\}$, we consider the semilinear elliptic problem
\begin{equation}\label{modelsemil}
\Delta_z\Gamma + \Gamma_+^\gamma=0 \quad \mbox{in}\ \mathbb{R}^2,
\quad \{\Gamma>0\}= B_1(0),
\end{equation}
which posesses a classical radially symmetric solution of the form
\begin{equation}\label{defGamma}
\begin{aligned}
\Gamma(z)
  =
\left\{\begin{array}{ll}
\nu(|z|)&\mathrm{if~}|z|\leq1\\
\nu'(1)\log|z|&\mathrm{if~}|z|>1
\end{array}
\right.,\\
\end{aligned}
\end{equation}
 with $\nu$ being the unique positive radial ground state solution of
 \[
 \Delta_z
 \nu+\nu^\gamma =0 \quad \text{in} \quad B_1(0), \quad \nu=0 \quad \text{on} \quad \partial B_1(0).
 \]
We emphasise the sign condition $\nu'(1)<0$, which will be invoked in several later stages. 

In addition, for any $\ve,\,\mu >0$, the rescaled profile 
\begin{equation}\label{erescgamma}
\Gamma_{\ve\mu}(z)=\Gamma\left(\frac{z}{\ve\mu}\right)
\end{equation}
is of class $C^1(\overline{B_{\ve\mu}(0)})$ and solves the elliptic equation
\begin{equation}\label{rescaledequation123}
    \Delta_z\Gamma_{\ve\mu}(z)+\frac{1}{\ve^{2}\mu^2}\left(\Gamma_{\ve\mu}(z)\right)_+^\gamma=0
\quad  \mbox{in}\ \mathbb{R}^2.
\end{equation}
In the rest of this section, we focus on the region $|z|<\delta$ for a small $\delta>0$, and define
\begin{equation}\label{defofgammahat}
\widehat{\Gamma}_{\ve\mu}(z)=\Gamma\left(\frac{z}{\ve\mu}\right)-\nu'(1)|\log(\ve\mu)|,
\end{equation}
which satisfies \[-\Delta_z\widehat{\Gamma}_{\ve\mu}(z) \rightharpoonup c\delta_0  \quad  \mbox{as} \quad \ve\mu \to 0, \quad \mbox{with} \quad c=\int_{\R^2} \Gamma^{\gamma}_{+}.\]
Based on the model problem \eqref{rescaledequation123}, we construct an approximate solution of \eqref{oriz} via elliptic singular perturbation methods, with \eqref{defofgammahat} as the reference profile.

We obtain the following Proposition.

\begin{prop}\label{propapprox}
Given $h>0$ and $P=(a,b)\in\R^2\setminus\{0\}$, let $R=\sqrt{a^2+b^2}$. For any $\mu>0$ and the variable $z$ in \eqref{var}, consider the  profile $\widehat{\Gamma}_{\ve\mu}(z)$ in \eqref{defofgammahat} and define the approximate solution of \eqref{oriz}
as \[\psi_{\mu}(z)=  \widehat{\Gamma}_{\ve\mu}(z)\left(1+c_1z_1+c_2|z|^2\right) +\frac{R^3}{2h(h^2+R^2)^{\frac{3}{2}}}H_{1}(z),\]
where \[c_1=\frac{Rh}{2(h^2+R^2)^{\frac{3}{2}}}, \quad c_2=\frac{3h^2R^2+R^4}{8(h^2+R^2)^{3}},\]
and $H_1$ solves  \[\Delta_z H_{1} +\frac{\operatorname{Re}(z^3)}{\ve^2\mu^2|z|^2}\left(\Gamma''\left(\frac{|z|}{\ve\mu}\right)-\frac{\Gamma'\left(\frac{|z|}{\ve\mu}\right)}{\frac{|z|}{\ve\mu}}\right)=0.\]
In the rescaled variable $y=\frac{z}{\ve\mu}$ and for any small $\delta>0$, in the region $|y| < \frac{\delta}{\ve\mu}$ we obtain
\[\ve^2\mu^2L_{x}(\psi_{\mu})= \Delta_y\Gamma(y) + \frac{3Rh^2+R^3}{2h(h^2+R^2)^{\frac{3}{2}}}\ve\mu y_1\left(\Gamma(y)\right)^{\gamma}_{+} + \ve^2\mu^2 E_{*},\]
where $E_{*}$ is a smooth function in $z=\ve\mu y$, uniformly bounded as $\ve\mu \to 0.$

In terms of the $x$ variable, the approximate stream function is given by
\begin{equation}\label{approxinx}
    \Psi_{\mu,P}(x)=\psi_{\mu}\left(A[P]^{-1}\left(x-P\right)\right),
\end{equation}
where $A[P]$ is defined in \eqref{var}.
\end{prop}
\begin{proof}
We first consider $\delta>0$ small and introduce the rescaled variable $y=\frac{z}{\ve\mu},$ so that for $|y|<\frac {\delta}{\ve\mu}$ we use \eqref{oper} to write
\begin{equation}\label{Bfirsterror}
\begin{aligned}
\ve^2\mu^2 B_{0}(\ve\mu y)[\widehat{\Gamma}_{\ve\mu}]
=& -\frac{2Rh}{(h^2+R^2)^{\frac 3 2}}\ve \mu y_1\partial_{y_1y_1}\widehat{\Gamma}_{\ve\mu}-\frac{2R}{h\sqrt{h^2+R^2}}\ve \mu y_2\partial_{y_1y_2}\widehat{\Gamma}_{\ve\mu}\\
&- \frac{R}{h\sqrt{h^2+R^2}}\left(1+\frac{2h^2}{h^2+R^2}\right)\ve\mu
\partial_{y_1}\widehat{\Gamma}_{\ve\mu}
+\ve^2\mu^2 E_1,
\end{aligned}
\end{equation}
where $E_1$ is a smooth function in $z=\ve \mu y$ and uniformly bounded as  $\ve\mu\to 0$.

In addition, direct calculations give
\begin{equation}\label{radialders}
\partial_{y_1} \widehat\Gamma_{\ve\mu}=\Gamma' \frac{y_1}{|y|},  \quad
\partial_{y_1y_1}\widehat\Gamma_{\ve\mu}
= \left(\Gamma''-\frac{\Gamma'}{|y|}\right)\frac{y_1^2}{|y|^2}+\frac{\Gamma'}{|y|}, \quad \partial_{y_1 y_2}\widehat\Gamma_{\ve\mu}=\left(\Gamma''-\frac{\Gamma'}{|y|}\right)\frac{y_1 y_2}{|y|^2},
\end{equation}
where $\Gamma'$ and $\Gamma''$ denote radial derivatives of $\Gamma(|y|)=\Gamma_{\ve\mu}(z)$ in \eqref{erescgamma}.

We then turn to complex variables by identifying $y=(y_1,y_2)$ with the complex number $y=y_1+iy_2$. Substituting \eqref{radialders} into \eqref{Bfirsterror} and using the identities
\[y_1y_2^2=\frac{y_1|y|^2}{4}-\frac{\operatorname{Re}(y^3)}{4}, \quad y_1^3=\frac{3y_1|y|^2}{4}+\frac{\operatorname{Re}(y^3)}{4},\] 
we arrive at
\begin{equation}\label{bgamma2}
\begin{aligned}
\ve^2\mu^2 B_{0}(\ve\mu y)[\widehat\Gamma_{\ve\mu}]=&-\frac{Rh}{(h^2+R^2)^{\frac{3}{2}}}\frac{\Gamma'\ve\mu y_1}{|y|}-\frac{R^3+4Rh^2}{2h(h^2+R^2)^{\frac{3}{2}}}\left(\Gamma''+\frac{\Gamma'}{|y|}\right)\ve\mu y_1 \\
&+\frac{\ve\mu R^3}{2h(h^2+R^2)^{\frac{3}{2}}}\left(\Gamma''-\frac{\Gamma'}{|y|}\right) \frac{\operatorname{Re}(y^3)}{|y|^2} +\ve^2\mu^2 E_1,
\end{aligned}
\end{equation}
with $E_1$ as in \eqref{Bfirsterror}. 

To eliminate the first error term in \eqref{bgamma2}, we modify the initial approximation $\widehat{\Gamma}_{\ve\mu}$ by setting \begin{equation}\label{psi1}
\psi_{1}(\ve\mu y)=(1+c_1\ve\mu y_1)\widehat{\Gamma}_{\ve\mu}, \quad c_1=\frac{1}{2}\frac{Rh}{(h^2+R^2)^{\frac 3 2}},
\end{equation}
where we obtain
\begin{align*}
\ve^2\mu^2 L_{x}\left(\psi_{1}\right)
&=\Delta_{y}\Gamma
-\left(\frac{3Rh^2+R^3}{2h(h^2+R^2)^{\frac{3}{2}}}\right)
 \left(\Gamma'' +\frac{\Gamma'}{|y|}\right)\ve\mu y_1
+ \frac{\ve\mu R^3}{2h(h^2+R^2)^{\frac{3}{2}}}
\left(\Gamma'' -\frac{\Gamma'}{|y|}\right)\frac{\operatorname{Re}(y^3) }{ |y|^2}\nonumber \\
& \hspace{4mm} + \ve^2\mu^2 B_{0}(\ve\mu y)\left[c_1\ve\mu y_1\widehat{\Gamma}_{\ve\mu}\right]+\ve^2\mu^2 E_1.
\end{align*} 
In addition, due to the expression
\[\ve^2\mu^2
B_0(\ve\mu y)[c_1\ve\mu y_1 \widehat\Gamma_{\ve\mu}]= -\frac{c_1\ve^2\mu^2 R}{h\sqrt{h^2+R^2}}\left(1+\frac{2h^2}{h^2+R^2}\right)\widehat{\Gamma}_{\ve\mu}+\ve^2\mu^2 E_2,\]
where $E_{2}$ is another smooth function in $\ve \mu y$ which is uniformly bounded as $\ve\mu \to 0$, we infer that the profile in \eqref{psi1} satisfies
\begin{equation}\label{jae}
\begin{aligned}
\ve^2\mu^2 L_{x}\left(\psi_{1}\right)
&=  \Delta_{y}\Gamma
-\frac{3Rh^2+R^3}{2h(h^2+R^2)^{\frac{3}{2}}}\left(\Gamma'' +\frac{\Gamma'}{|y|}\right)\ve\mu y_1
-
 \frac{\ve^2\mu^2 (3h^2R^2+R^4)}{2(h^2+R^2)^3} \widehat{\Gamma}_{\ve\mu}
\\
& \hspace{4mm} + \frac{\ve\mu R^3}{2h(h^2+R^2)^{\frac{3}{2}}}
\left(\Gamma^{''} -\frac{\Gamma'}{|y|}\right)\frac{\operatorname{Re}(y^3) }{|y|^2}
+\ve^2\mu^2(E_1+E_2).
\end{aligned}
\end{equation}
To remove the error term proportional to $\widehat{\Gamma}_{\ve\mu}$ in \eqref{jae}, we first note that  
\[
   \Delta_{y}\left(c_2 \ve^2\mu^2|y|^2\widehat\Gamma_{\ve\mu}\right)
=c_2\ve^2\mu^2\left(4\widehat{\Gamma}_{\ve\mu}+4 y\cdot\nabla_{y}\Gamma+|y|^2\Delta_{y}\Gamma\right),\]
and then adjust the approximation to 
\begin{equation*}
\psi_{2}(\ve\mu y)=\left(1+c_1\ve\mu y_1+c_2\ve^2\mu^2|y|^2\right)\widehat\Gamma_{\ve\mu}, \quad c_2= \frac{3h^2R^2+R^4}{8(h^2+R^2)^3},
\end{equation*}
for which
\begin{equation}\label{errorsecondapprox}
\begin{aligned}
\ve^2 \mu^2 L_{x}\left(\psi_{2}\right) &=  \Delta_{y}\Gamma
-\frac{3Rh^2+R^3}{2h(h^2+R^2)^{\frac 3 2}}
 \left(\Gamma'' +\frac{\Gamma'}{|y|}\right)\ve\mu y_1
  + \frac{\ve\mu R^3}{2h(h^2+R^2)^{\frac 3 2}} \left(\Gamma'' -\frac{\Gamma'}{|y|}\right)\frac{\operatorname{Re}(y^3)}{|y|^2}\\
  &+\ve^2\mu^2(E_1 +E_2+E_3),
\end{aligned}
\end{equation}
with $E_3$ having the same properties as $E_1$ and $E_2$ above.

To cancel the cubic term in \eqref{errorsecondapprox}, we now introduce the radial correction
\[
h_{1}(s)
= s^3\int_s^1\frac{\rm dr}{r^7}
\int_0^r \frac{t^5}{\ve^2\mu^2}\left(\Gamma''\left(\frac{|t|}{\ve\mu}\right) -\frac{\Gamma'\left(\frac{|t|}{\ve\mu}\right)}{\frac{|t|}{\ve\mu}}\right) \rm d t,
\]
which solves
\[h''_{1}+\frac{1}{s} h'_{1}-\frac{9}{s^2}h_{1} +\frac{s}{\ve^2\mu^2}\left(\Gamma''\left(\frac{|s|}{\ve\mu}\right)-\frac{\Gamma'\left(\frac{|s|}{\ve\mu}\right)}{\frac{|s|}{\ve\mu}}\right)=0.\]
We emphasise that the term $\frac{s}{\ve^2\mu^2}\left(\Gamma''\left(\frac{|s|}{\ve\mu}\right)-\frac{\Gamma'\left(\frac{|s|}{\ve\mu}\right)}{\frac{|s|}{\ve\mu}}\right)$ is uniformly bounded as $\ve\mu\to 0$.
Indeed, by \eqref{defGamma} we get $u\left(\Gamma''(u)-\frac{\Gamma'(u)}{u}\right)=\mathcal{O}\left(\frac{1}{u}\right)$ as $u \to \infty,$ while $\lim\limits_{u\to 0}\frac{\Gamma'(u)}{u}=\Gamma''(0)$ implies $\left(\Gamma''(u)-\frac{\Gamma'(u)}{u}\right)=\mathcal{O}(u)$ as $u \to 0.$ In particular, $h_{1}$ is a smooth function that remains uniformly bounded as $\ve\mu \to 0,$ satisfying $h_{1}(s)=\mathcal{O}(s^3)$ as $s\to 0$. 

Writing $y=|y|e^{i\theta}$, we set $H_{1}(z)=h_{1}(|z|)\cos(3\theta)$ to obtain a solution of
\begin{equation*}
\Delta_{y} H_{1}+\ve\mu\left(\Gamma''(y) -\frac{\Gamma'(y)}{|y|}\right)\frac{\operatorname{Re}(y^3)}{|y|^2} =0.
\end{equation*}
This leads us to the final approximation 
\begin{equation}\label{psi3}
\psi_{\mu}(\ve\mu y)=\left(1+c_1\ve \mu y_1+c_2\ve^2\mu^2|y|^2\right)\widehat\Gamma_{\ve\mu}
+\frac{R^3}{2h(h^2+R^2)^{\frac{3}{2}}}H_{1}(\ve\mu y),
\end{equation}
so that \eqref{modelsemil} yields
\begin{equation*}
\begin{aligned}
\ve^2 \mu^2 L_{x}\left(\psi_{\mu}\right)= &\Delta_y\Gamma
+\frac{3Rh^2+R^3}{2h(h^2+R^2)^\frac{3}{2}}
 \ve\mu y_1 \Gamma_{+}^{\gamma}
 +\ve^2\mu^2 E_{*},
\end{aligned}
\end{equation*}
with $E_{*}$ an additional smooth function in $\ve\mu y$, uniformly bounded as $\ve\mu \to 0$. 

To this end, we use \eqref{var} to express the approximate stream function \eqref{psi3} around $P=(a,b)$ in the original coordinates $x=(x_1,x_2)$ as
\begin{equation*}
\Psi_{\mu,P}(x)=\psi_{\mu}\left(A[P]^{-1}(x-P)\right).
\end{equation*}
\end{proof}
\section{Construction of a global multi-vortex approximation}\label{section3}
Given an integer $N\geq 1$, in this section we construct an approximate stream function for an $N$-vortex solution of \eqref{eqtosolve}. The approximation is designed to resemble a superposition of Dirac masses concentrated at $N$ distinct points in the plane, whose pairwise distance shrinks as $\ve\to 0.$ As will be shown below, it is obtained by combining suitably rescaled copies of the building block introduced in \eqref{approxinx}, centered at each vortex point. 

More specifically, we let $P_j=(a_j,b_j),\,j=1,\dots,N$ be as in 
\eqref{formpoints}, consider positive scaling parameters $\mu_j>0$ to be determined, and employ \eqref{approxinx} to define
\begin{equation}\label{profj}
\Psi_j(x)=\Psi_{\mu_j,P_j}(x).
\end{equation}
 Using Proposition \ref{propapprox}, for any $\delta>0$ small and the variable $z=A[P_j]^{-1}(x-P_j)=\ve\mu_j y$  in \eqref{var}, for each $j\in\{1,\dots,N\}$ we infer that
\begin{equation}\label{erroraroundeachj}
\ve^2\mu_j^2L_{x}(\Psi_j)=\Delta_{y}\Gamma +\frac{3R_jh^2+R_j^3}{2h(h^2+R_j^2)^{\frac 3 2}}\ve \mu_j y_1 \Gamma^{\gamma}_{+} +\ve^2\mu_j^2 E_{*,j}, \quad |y|<\frac{\delta}{\ve\mu_j},
\end{equation}
where $R_j=\sqrt{a_j^2+b_j^2}$ and $E_{*,j}$ is a smooth function in $\ve\mu_j y,$ which is uniformly bounded as $\ve\mu_j \to 0.$ 

Nevertheless, since the points $P_j$ cluster near $(r_0,0)$ on a scale $\mathcal{O}(|\log\ve|^{-1})$, we multiply each $\Psi_j$ by a smooth cut-off function $\eta_0:\R^2\to[0,1]$ 
to obtain the profile
\[\eta_{0}(x)\sum_{j=1}^{N}\kappa_j\Psi_j(x),\] 
where 
\begin{equation*}
    \eta_{0}(x)=\eta\left(|x-(r_0,0)|\right),
\end{equation*}
with $\eta:\R\to [0,1]$ denoting a smooth cut-off satisfying 
\begin{equation}\label{defeta}
\eta(s)=1 \quad \mbox{for} \quad s\leq \frac 1 2, \quad \eta(s)=0 \quad \mbox{for}\quad s \geq 1.
\end{equation}
In this way, the approximate stream function is extended to the whole $\R^2$, while setting $z_j=([z_j]_1,[z_j]_2)=A[P_j]^{-1}(x-P_j)$ and using \eqref{erroraroundeachj}, a direct computation yields 
\[L_{x}\left(\eta_0\sum_{j=1}^{N}\kappa_j\Psi_j\right)=\eta_{0}\sum_{j=1}^{N}\kappa_j\left(\Delta_{z_j}\Gamma_{\ve\mu_j}+\frac{3R_j h^2+R_j^3}{2h(h^2+R_j^2)^{\frac 3 2}}\frac{[z_j]_1}{\ve^2\mu_j^2}(\Gamma_{\ve\mu_j})^{\gamma}_{+}\right) +g(x),
\]
where 
\[g(x)=\eta_{0}\sum_{j=1}^{N}\kappa_j E_{*,j}+\sum_{j=1}^{N}\kappa_j\left[L_{x}\left(\eta_{0}\Psi_j\right)-\eta_{0}L_{x}\left(\Psi_j\right)\right].
\]
One can verify that the function $g(x)$ is compactly supported, with  $\|g\|_{L^{\infty}\left(\R^2\right)} \leq C$ for some $C>0.$ 

To proceed, we refine the preliminary ansatz  $\eta_{0}\sum\limits_{j=1}^{N}\kappa_j\Psi_j$ by adding a global term that eliminates the bounded error term $g(x)$, resulting in a more accurate approximate solution. 

In particular, we let $H_{2\ve}(x)$ be a solution of
\begin{equation*}
L_{x}(H_{2\ve})+g(x)=0 \quad \text{in} \quad \R^2,
\end{equation*}
whose existence follows from the linear theory of Proposition \ref{propouter}, and which satisfies the growth estimate
\begin{equation}\label{boundh2e}
|H_{2\ve}(x)|\leq C\left(1+|x|^2\right)\quad \forall x\in\R^2.
\end{equation}
In addition, we fix the additive constant by imposing the condition
\begin{equation*}
    H_{2\varepsilon}(r_0,0)=0.
\end{equation*}
In terms of the variable $z_j=A[P_j]^{-1}(x-P_j),\,z_j=([z_j]_1,[z_j]_2),$ we define the final approximate stream function 
\begin{equation}\label{globalapprox}
    \Psi_0(x)=\eta_{0}(x)\sum_{j=1}^{N}\kappa_j\left[\widehat{\Gamma}_{\ve\mu_j}(z_j)\left(1+c_{1,j}[z_j]_1+c_{2,j}|z_j|^2\right)+\frac{R_j^3}{2h(h^2+R_j^2)^{\frac 3 2}}H_{1}(z_j)\right] +H_{2\ve}(x),
\end{equation}
where $\widehat\Gamma_{\ve\mu_j}$ is given by \eqref{defofgammahat}, while the function $H_{1}$ and the constants $c_{1,j}$,\,$c_{2,j}$, \,$j=1,\dots,N$ are as in Proposition \ref{propapprox}, with $P=(a,b)$ replaced by $P_j=(a_j,b_j)$. 

Collecting the above, we conclude that the function $\Psi_0$ in \eqref{globalapprox} satisfies 
\begin{equation}\label{lxofglobal}
    L_{x}(\Psi_0)= \eta_0\sum_{j=1}^{N}\kappa_j\left(\Delta_{z_j} \Gamma_{\ve\mu_j}+\frac{3R_jh^2+R_j^3}{2h(h^2+R_j^2)^{\frac 3 2}}\frac{[z_j]_1}{\ve^2\mu_j^2}\left(\Gamma_{\ve\mu_j}\right)^{\gamma}_{+}\right), \quad z_j=A[P_j]^{-1}(x-P_j),
\end{equation}
which in view of \eqref{defGamma}-\eqref{rescaledequation123}, the result in \eqref{lxofglobal} further yields
\begin{equation}\label{Lxsupport}
\operatorname{supp}
L_x(\Psi_0)\subset\bigcup_{j=1}^N \{x\in\mathbb{R}^2:\ |A[P_j]^{-1}(x-P_j)|<\ve\mu_j\}.
\end{equation}
 It is worth recalling that $\eta_0$ denotes the cut-off in \eqref{defeta} and the notation $R_j=\sqrt{a_j^2+b_j^2}$, while we also remark that the profile $\Psi_0$ in \eqref{globalapprox} is smooth and globally defined in $\R^2$, with its dependence on the points $P_1,\dots,P_N$ and the positive scaling parameters $\mu_1,\dots,\mu_N$ determined by \eqref{profj}.
\section{Properties of the scaling parameters}\label{Section 4}
The aim of this section is to choose the scaling parameters $\mu_j,\, j=1,\dots,N$ introduced in \eqref{profj}, and to describe how they depend on the vortex points $P_j$ in \eqref{formpoints}. We use the decomposition
\begin{equation}\label{scalingansatz}
\mu_j=\mu_j^{0}+\mu^{*}_j, \qquad j=1,\dots,N,
\end{equation}
where $\mu_j^{0}$ designates the leading order scaling parameter around $P_j$ and $\mu_j^{*}$ is a small correction.

For each $j\in\{1,\dots,N\}$, the parameter $\mu_j^{0}$ will be chosen so that
\[
\log(\mu^{0}_{j})=\mathcal{O}\left(\log|\log\ve|\right)\quad \text{as }\ve\to 0,
\]
while for the perturbation $\mu_j^{*}$ we assume a priori that it satisfies
\begin{equation}\label{aprioriscalingperp}
\log\bigg(1+\frac{\mu_{j}^{*}}{\mu_{j}^{0}}\bigg)
=\mathcal{O}\left(\ve^{1+\tilde{\sigma}}\right), \quad \mbox{for a small} \quad \tilde{\sigma}>0.
\end{equation}
Although the explicit choice of $\mu_j^{*}$ and the precise value of $\tilde{\sigma}>0$ in \eqref{aprioriscalingperp} are deferred to
Proposition \ref{propsec8}, we stress that these corrections are vital for carrying out the construction, as discussed in Remark \ref{decoupling}. 

Currently, our objective is to explain in more detail the role of the leading order parameters $\mu_j^{0}$. They are chosen so as to partially eliminate the radial terms of size
$\mathcal{O}(\log|\log\ve|)$ arising in the local expansion of
$\Psi_{0}-\frac{\alpha}{2}|\log\ve|\,|x|^{2}$
near each vortex point $P_j$.

More concretely, for $d>0$ as in \eqref{unifdistance}, we fix $\delta>0$ sufficiently small so that
\begin{equation}\label{cond1delta}
\delta < \frac{\sqrt{h^2+r_0^2}}{4h}d,
\end{equation}
and consider the region
\begin{equation}\label{firstdelta}
|A_{i}^{-1}\left(x-P_i\right)|< \frac{\delta}{|\log\ve|},
\end{equation}
where $\eta_{0}(x)= 1$ due to \eqref{defeta}.  

We have the following Proposition.
\begin{prop}\label{propscaling}
For each $j=1,\dots,N$, let $A_j=A[P_j]$ as in \eqref{var}, evaluated at the point $P_j$ in \eqref{formpoints}.  Fix $i\in\{1,\dots,N\},$ and let $\mu_i=\mu_i^0 +\mu_i^*$ satisfy
\begin{equation}\label{scalingmain}
\begin{aligned}
-\kappa_i\nu'(1)\log\left(\mu_{i}^{0}\right)&=\sum_{j\neq i} \kappa_j\nu'(1)\log\left(|A_j^{-1}\left(P_i-P_j\right)|\right)\bigg[1+c_{1,j}\left[A_j^{-1}\left(P_i-P_j\right)\right]_{1}+c_{2,j}\left|A_j^{-1}\left(P_i-P_j\right)\right|^2\bigg]\\
&\hspace{4mm}+H_{2\ve}(P_i),
\end{aligned}
\end{equation}
and
\begin{equation}\label{scalingpert}
-\kappa_i\nu'(1)\log\left(1+\frac{\mu_i^{*}}{\mu_i^{0}}\right)=\mathcal{O}\left(\ve^{1+\tilde\sigma}\right)\quad \mbox{for some} \, \tilde\sigma>0,
\end{equation}
with 
\begin{equation}\label{constantsj}
    c_{1,j}=\frac{R_j h}{2(h^2+R_j^2)^{\frac 3 2}},\quad c_{2,j}=\frac{3h^2R_j^2+R_j^4}{8(h^2+R_j^2)^3}, \quad R_j=\sqrt{a_j^2+b_j^2}, \quad j=1,\dots,N.
\end{equation}
For $\delta>0$ as in \eqref{cond1delta} and the variable \[z=A_i^{-1}(x-P_i),\, z=\ve\mu_i y,\] in the region $|z|<\frac{\delta}{|\log\ve|}$ the approximate stream function $\Psi_{0}$ in \eqref{globalapprox} has the expansion
\begin{equation}\label{exppsiminusfinal}
\begin{aligned}
&\frac{1}{\kappa_{i}}\left(\Psi_{0}-\frac{\alpha}{2}|\log\ve||x|^2\right)=\\
&\left(1+c_{1,i}\ve\mu_{i}y_1+c_{2,i}\ve^2\mu_{i}^2|y|^2\right)\Gamma(y) -\frac{\alpha}{2\kappa_{i}}|\log\ve||P_i|^2-\nu'(1)|\log\ve|+\nu'(1)\log\left(1+\frac{\mu_{i}^{*}}{\mu_{i}^{0}}\right)\\
&+\ve \mu_{i}y_1\Bigg[|\log\ve|\left(-c_{1,i}\nu'(1)-\frac{\alpha R_i h}{\kappa_i\sqrt{h^2+R_i^2}}\right)+\nu'(1)c_{1,i}\log(\mu_{i})\\
&\hspace{14mm}+\sum_{j\neq i}\frac{\kappa_j}{\kappa_i}\nu'(1)\frac{[A_j^{-1}(P_i-P_j)]_1}{|A_j^{-1}(P_i-P_j)|^2}\left(1+c_{1,j}[A_j^{-1}(P_i-P_j)]_1+c_{2,j}|A_{j}^{-1}(P_i-P_j)|^2\right)\\
&\hspace{14mm}+\sum_{j\neq i}\frac{\kappa_j}{\kappa_{i}}\nu'(1)\log\left(|A_j^{-1}\left(P_i-P_j\right)|\right)\left(c_{1,j}+2c_{2,j}[A_j^{-1}(P_i-P_j)]_1\right)+\frac{1}{\kappa_i}\left(A_{i}(1,0)^{T}\right)\nabla H_{2\ve}(P_i)\Bigg]\\
&+\ve \mu_i y_2 \Bigg[\sum_{j \neq i}\frac{\kappa_j}{\kappa_i}\nu'(1)\frac{[A_j^{-1}(P_i-P_j)]_{2}}{|A_j^{-1}(P_i-P_j)|^2}\left(1+c_{1,j}[A_j^{-1}(P_i-P_j)]_{1}+c_{2,j}|A_j^{-1}(P_i-P_j)|^2\right)\\
&\hspace{14mm}+\sum_{j\neq i}\frac{2\kappa_j}{\kappa_{i}}\nu'(1)\log\left(|A_j^{-1}\left(P_i-P_j\right)|\right)c_{2,j}[A_j^{-1}(P_i-P_j)]_2+\frac{1}{\kappa_i}\left(A_i(0,1)^{T}\right)\nabla H_{2\ve}(P_i)\Bigg]\\
&+\mathcal{O}\left(\ve^2 \mu_i ^2 |\log\ve|^2 |y|^2 \right) \quad \mbox{as} \quad \ve \to 0.
\end{aligned}
\end{equation} 
\end{prop}
\begin{proof}
 To begin with, we fix $i\in\{1,\dots,N\}$ and let $A_{i}=A[P_{i}]$ as in \eqref{var}. Since the points in \eqref{formpoints} are of the form \[P_i=(r_0+\tilde{r},0)+\frac{1}{|\log\ve|}\widehat{P_i}, \quad P_i=(a_i,b_i),\quad \widehat{P_i}=(\hat{a}_i,\hat{b}_i),\] we obtain the matrix expansions 
 \begin{equation}\label{expinverse}
 \begin{aligned}
 A_{i}^{-1}&=\begin{pmatrix}
     \frac{a_{i}\sqrt{h^2+R_i^2}}{R_i h} & \frac{b_i\sqrt{h^2+R_i^2}}{R_i h}\\
     \frac{-b_i}{R_i} & \frac{a_i}{R_i}\end{pmatrix}=\begin{pmatrix}
         \frac{\sqrt{h^2+r_0^2}}{h} & 0 \\ 0 & 1 \end{pmatrix} + \frac{\log|\log\ve|}{|\log\ve|}\widehat{A}_{i},\\ 
A_j^{-1}A_i&= \begin{pmatrix} 1 & 0 \\ 0 & 1 \end{pmatrix} +\frac{\log|\log\ve|}{|\log\ve|}\widehat{I}_{ij}, \quad \mbox{for all}\, i \neq j,
\end{aligned}
\end{equation}
where $\widehat{A}_{i}$ and $\widehat{I}_{ij}$ are $2\times 2$ matrices with smooth entries depending on $(\tilde{r},\hat{a}_i,\hat{b}_i)$ and $(\tilde{r},\hat{a}_i,\hat{b}_i,\hat{a}_j,\hat{b}_j),$ respectively, which are uniformly bounded as $\ve\to 0.$

In addition, in the region $|z|<\frac{\delta}{|\log\ve|}$ and the rescaled variable $y=\frac{z}{\ve\mu_i}$, as $\ve\to 0$ we get 
\begin{equation}\label{expanH}
\begin{aligned}
    H_{2\ve}(x)&=H_{2\ve}(P_i)+\ve \mu_i\left(A_i y\right)\cdot\nabla H_{2\ve}(P_i) +\mathcal{O}(\ve^2 \mu_i^2 |y|^2),\\
    H_1(\ve \mu_i y)&=\mathcal{O}(\ve^3 \mu_i^3 |y|^3),\\
\end{aligned}
\end{equation}
as well as
\begin{equation}\label{expansionGamma}
\begin{aligned}
&\widehat{\Gamma}\left(\frac{A_j^{-1}A_iz+A_{j}^{-1}(P_i-P_j)}{\ve\mu_j}\right)=\nu'(1)\log\left(\frac{|A_j^{-1}A_i\ve\mu_i y +A_j^{-1}(P_i-P_j)|}{\ve\mu_j}\right)-\nu'(1)|\log(\ve\mu_j)|\\
&=\nu'(1)\log(|A_j^{-1}(P_i-P_j)|)+\nu'(1)\frac{A_j^{-1}(P_i-P_j)}{|A_j^{-1}(P_i-P_j)|^2}\cdot \ve\mu_i A_j^{-1}A_iy +\mathcal{O}\left(\frac{|\ve\mu_i A_j^{-1}A_iy|^2}{|A_j^{-1}(P_i-P_j)|^2}\right).
\end{aligned}
\end{equation}
Using the definition of $\Psi_0$ in \eqref{globalapprox}, we write
\begin{equation}\label{psiminusquad}
\begin{aligned}
&\Psi_0 -\frac{\alpha}{2}|\log\ve||x|^2 \\
&=\kappa_i(1+c_{1,i}\ve\mu_iy_1+c_{2,i}\ve^2\mu_i^2|y|^2)\Gamma(y)\\
&\hspace{3mm}+\kappa_i\left(-\nu'(1)|\log\ve|+\nu'(1)\log(\mu^{0}_{i})+\nu'(1)\log\Big(1+\frac{\mu^{*}_i}{\mu^{0}_i}\Big)\right)(1+c_{1,i}\ve\mu_iy_1+c_{2,i}\ve^2\mu_i^2|y|^2)\\
&\hspace{3mm}+\kappa_i\frac{R_i^3}{2h(h^2+R_i^2)^{\frac 3 2}}H_{1}\left(\ve \mu_i y\right) +H_{2\ve}(P_i+A_i\ve\mu_i y)\\
&\hspace{3mm}-\frac{\alpha}{2}|\log\ve||P_i|^2-\frac{\alpha R_ih|\log\ve|}{\sqrt{h^2+R_i^2}}\ve\mu_i y_1-\frac{\alpha}{2}|\log\ve|\ve^2\mu_i^2|A_iy|^2\\
&\hspace{3mm}+\sum_{j\neq i}\kappa_j\left(\nu'(1)\log\left(\frac{|A_j^{-1}A_i\ve \mu_i y +A_j^{-1}(P_i-P_j)|}{\ve \mu_j}\right)-\nu'(1)|\log(\ve\mu_j)|\right)\\
&\hspace{7mm}\times\left(1+c_{1,j}[A_j^{-1}A_i\ve\mu_iy+A_j^{-1}(P_i-P_j)]_1+c_{2,j}|A_j^{-1}A_i\ve\mu_iy+A_j^{-1}(P_i-P_j)|^2\right)\\
&\hspace{3mm}+\sum_{j \neq i}\kappa_j\frac{R_j^3}{2h(h^2+R_j^2)^{\frac 3 2}}H_1\left(A_j^{-1}A_i\ve\mu_i y +A_j^{-1}(P_i-P_j)\right).
\end{aligned}
\end{equation}
Substituting \eqref{expanH} and \eqref{expansionGamma} into \eqref{psiminusquad} and then using  \eqref{expinverse}, we arrive at the expansion in \eqref{exppsiminusfinal}, where constant terms are removed by the choice of $\mu_i^{0}$ in \eqref{scalingmain} and the contribution of $\mu_i^{*}$ is controlled by \eqref{scalingpert}. This concludes the proof.
\end{proof}
\begin{remark}
Combining \eqref{formpoints} with the relation \eqref{scalingmain}, we verify that for each $i=1,\dots,N$, the leading order scaling parameters $\mu_i^{0}$ in Proposition \ref{propscaling} satisfy
    \[-\nu'(1)\log(\mu_i^{0})=\mathcal{O}\left(\log|\log\ve|\right) \quad \mbox{as} \quad \ve \to 0.\]
    In particular, the corresponding vortex core scales satisfy $\ve \mu_i=\mathcal{O}(\ve|\log\ve|^2)$, and are therefore asymptotically much smaller than the pairwise distance $\mathcal{O}(|\log\ve|^{-1})$ between the vortex centers.
 For later use, we set 
 \begin{equation}\label{maxscaling}
     \mu=\max_{i=1,\dots,N}\mu_i.
 \end{equation}
\end{remark}
\section{Choice of the vorticity profile and approximation error estimates}\label{Section 5}
In the present section, the objective is to select a suitable nonlinearity $F$ in \eqref{eqtosolve} so that the corresponding vorticity profile \[W(x)=F\left(\Psi-\frac{\alpha}{2}|\log\ve||x|^2\right)\] satisfies \eqref{propvorticity}, and to quantify the approximation error arising from the substitution of $\Psi_0(x)$ from \eqref{globalapprox} into \eqref{eqtosolve}. 

In this regard, we introduce the error operator
\begin{equation}\label{definitionerroroper}
S(\Psi)(x)=L_{x}(\Psi)+F\left(\Psi-\frac{\alpha}{2}|\log\ve||x|^2\right), \quad x\in\R^2,
\end{equation}
so that $\Psi$ is an exact solution if and only if $S(\Psi)=0$.

Inspired by the structure of \eqref{rescaledequation123}, for $\gamma >3$ and $s_{+}\defeq \max(s,0)$ we define the nonlinearity 
\begin{equation}\label{nonlinearity}
F(s)=\sum_{i=1}^{N}\frac{\kappa_i}{\ve^2\mu_i^2}F_i(s), \quad
F_i(s)=\left(\frac{s}{\kappa_i}+\nu'(1)|\log\ve|+\frac{\alpha}{2\kappa_i}|\log\ve||P_i|^2\right)^{\gamma}_{+}\eta_{i}\left(\frac{s}{\kappa_i}\right),
\end{equation}
where $\eta_{i}:\R\to [0,1],\, i=1,\dots,N,$ are smooth cut-off functions specified as follows.

Consider the inner region $|A_i^{-1}(x-P_i)|<\frac{\delta}{|\log\ve|}$ centered at the vortex point $P_i$. In the rescaled variable $y=\frac{A_i^{-1}(x-P_i)}{\ve\mu_i}$, the boundary of this region corresponds to $|y|=\frac{\delta}{\ve\mu_i|\log\ve|},$ where the expansion derived in Proposition \ref{propscaling} gives
\[\frac{1}{\kappa_i}\left(\Psi_0-\frac{\alpha}{2}|\log\ve||x|^2\right)=\nu'(1)\log(\delta)\left(1+\tilde{s}(\delta)\right)-\nu'(1)\log|\log\ve|-\nu'(1)\log(\mu_i^{0})-\frac{\alpha}{2\kappa_i}|\log\ve||P_i|^2+o(1),\]
with $\tilde{s}(\delta)$ being a smooth bounded function depending on $\delta$, and $o(1) \to 0$ as $\ve \to 0.$

With this in mind, we define the cut-off $\eta_{i}$ as
\begin{equation}\label{cutoffj}
\eta_{i}(s)=\begin{cases}
    1, \quad \mbox{for}\quad s \geq -\nu'(1)\log|\log\ve|-\nu'(1)\log(\mu_i^{0})-\frac{\alpha}{2\kappa_i}|\log\ve||P_i|^2+2\ell_{i,\ve},\\
    0, \quad \mbox{for}\quad s \leq -\nu'(1)\log|\log\ve|-\nu'(1)\log(\mu_i^{0})-\frac{\alpha}{2\kappa_i}|\log\ve||P_i|^2+\ell_{i,\ve},
\end{cases}
\end{equation}
for appropriately chosen $\ell_{i,\ve}=\nu'(1)\log\left(\delta\right)\left(1+\tilde{s}\left(\delta\right)\right)+o(1)$, so that
\begin{equation}\label{rangecutoffi}
\begin{cases}
\eta_{i}\left(\frac{1}{\kappa_{i}}\left(\Psi_0-\frac{\alpha}{2}|\log\ve||x|^2\right)\right)=1 \quad \mbox{for} \quad |A_{i}^{-1}(x-P_i)|\leq \frac{\delta^2}{|\log\ve|},\\
\eta_{i}\left(\frac{1}{\kappa_{i}}\left(\Psi_0-\frac{\alpha}{2}|\log\ve||x|^2\right)\right)=0 \quad \mbox{for} \quad |A_{i}^{-1}(x-P_i)|\geq \frac{\delta}{|\log\ve|}.
\end{cases}
\end{equation}
In the above, $\delta>0$ is independent of $\ve>0$ as in \eqref{firstdelta} and \eqref{cond1delta}, and may be taken smaller if necessary.

The following Proposition holds.
\begin{prop}\label{errorprop1}
    Consider $r_0>0,\, h>0$, and let $\alpha$ be the fixed rotational speed in \eqref{uniformspeed}. For the approximate stream function $\Psi_{0}$ in \eqref{globalapprox}, the nonlinear function $F$ in \eqref{nonlinearity}, and $\delta>0$ sufficiently small satisfying \eqref{cond1delta}, there exists a constant $C>0$ such that for all $\ve >0$ small and each $i\in\{1,\dots,N\}$, the error function in \eqref{definitionerroroper} satisfies
    \[\ve^2\mu_i^2\left| S\left(\Psi_{0}\right)\right|\leq C\ve\mu_i|\log\ve|\left(\Gamma\left(\frac{y}{2}\right)\right)^{\gamma}_{+},\]
where $z=A_i^{-1}(x-P_i), \,y=\frac{z}{\ve\mu_i},$ and $\operatorname{supp} \left(\Gamma(y)\right)^{\gamma}_{+} \subset B_1 (0).$
\end{prop}
\begin{proof}
To estimate the approximation error 
\begin{equation}\label{rescalederroroper}
\begin{aligned}
    &\ve^2\mu_i^2 S(\Psi_{0})=\ve^2\mu_i^2 L_{x}(\Psi_{0})+\ve^2\mu_i^2\sum_{j=1}^{N}\frac{\kappa_j}{\ve^2\mu_j^2}F_{j}\left(\Psi_0-\frac{\alpha}{2}|\log\ve||x|^2\right),
\end{aligned}
\end{equation}
with $L_x(\Psi_0)$ and $F_j\left(\Psi_0-\frac{\alpha}{2}|\log\ve||x|^2\right)$ given in \eqref{lxofglobal} and \eqref{nonlinearity}, respectively, we consider five distinct regions in the variable  \[z=A_i^{-1}(x-P_i),\quad z=\ve\mu_i y\] as follows.
\\
\text{\bf Case 1:} $|y|< 1.$
\\
Due to \eqref{cond1delta} and \eqref{rangecutoffi}, in this region we have $\eta_{i}=1,$ and
$\eta_j=0$ for all $j\neq i$ as $|P_i-P_j|\geq \frac{\delta}{|\log\ve|}$ , thus the nonlinear component in \eqref{rescalederroroper} simplifies to 
\begin{equation}\label{approxerrfirst}
\begin{aligned}
&\ve^2\mu_i^2\sum_{j=1}^{N}\frac{\kappa_j}{\ve^2\mu_j^2}\left(\frac{1}{\kappa_j}\left(\Psi_0-\frac{\alpha}{2}|\log\ve||x|^2\right)+\nu'(1)|\log\ve|+\frac{\alpha}{2\kappa_j}|\log\ve||P_j|^2\right)^{\gamma}_{+}\eta_{j}\left(\frac{1}{\kappa_j}\left(\Psi_0-\frac{\alpha}{2}|\log\ve||x|^2\right)\right)\\
&=\kappa_i\Bigg[\Gamma(y)+(c_{1,i}\ve\mu_iy_1+c_{2,i}\ve^2\mu_i^2|y|^2)\Gamma(y)+\nu'(1)\log\left(1+\frac{\mu^{*}_i}{\mu^{0}_i}\right)+\ve\mu_i y_1\mathcal{A}_{1,i}(P)+\ve\mu_i y_2\mathcal{A}_{2,i}(P)\\
&\hspace{10mm}+\mathcal{O}\left(\ve^2\mu_i^2|\log\ve|^2|y|^2\right)\Bigg]^{\gamma}_{+},
\end{aligned}
\end{equation}
where for $P=(P_1,\dots,P_N)$ as in \eqref{formpoints}, we have adopted the notation
\begin{equation}\label{mode1func}
\begin{aligned}
\mathcal{A}_{1,i}(P)&=|\log\ve|\left(-c_{1,i}\nu'(1)-\frac{\alpha R_i h}{\kappa_i\sqrt{h^2+R_i^2}}\right)+\nu'(1)c_{1,i}\log(\mu_{i}^{0})+\nu'(1)c_{1,i}\log\left(1+\frac{\mu^*_{i}}{\mu^{0}_{i}}\right)\\
&\hspace{3mm}+\sum_{j\neq i}\frac{\kappa_j}{\kappa_{i}}\nu'(1)\log\left(|A_j^{-1}\left(P_i-P_j\right)|\right)\left(c_{1,j}+2c_{2,j}[A_j^{-1}(P_i-P_j)]_1\right)\\
&\hspace{3mm}+\sum_{j\neq i}\frac{\kappa_j}{\kappa_i}\nu'(1)\frac{[A_j^{-1}(P_i-P_j)]_1}{|A_j^{-1}(P_i-P_j)|^2}\left(1+c_{1,j}[A_j^{-1}(P_i-P_j)]_1+c_{2,j}|A_{j}^{-1}(P_i-P_j)|^2\right)\\
&\hspace{3mm}+\frac{1}{\kappa_i}\left(A_{i}(1,0)^{T}\right)\nabla H_{2\ve}(P_i),\\
\mathcal{A}_{2,i}(P)&=\sum_{j \neq i}\frac{2\kappa_j}{\kappa_i}\nu'(1)\log\left(|A_j^{-1}(P_i-P_j)|\right)c_{2,j}[A_j^{-1}(P_i-P_j)]_{2}\\
&\hspace{3mm}+\sum_{j \neq i}\frac{\kappa_j}{\kappa_i}\nu'(1)\frac{[A_j^{-1}(P_i-P_j)]_{2}}{|A_j^{-1}(P_i-P_j)|^2}\left(1+c_{1,j}[A_j^{-1}(P_i-P_j)]_{1}+c_{2,j}|A_j^{-1}(P_i-P_j)|^2\right)\\
&\hspace{3mm}+\frac{1}{\kappa_i}\left(A_i(0,1)^{T}\right)\nabla H_{2\ve}(P_i).
\end{aligned}
\end{equation}
In fact, a careful analysis of the terms in \eqref{mode1func} allows us to further write 
\begin{equation}\label{logmode1fun}
    \mathcal{A}_{1,i}(P)=|\log\ve|\,\widehat{\mathcal{A}}_{1,i}(P), \quad 
\mathcal{A}_{2,i}(P)=|\log\ve|\,\widehat{\mathcal{A}}_{2,i}(P),
\end{equation}
for some smooth functions $\widehat{\mathcal{A}}_{1,i}(P),\,\widehat{\mathcal{A}}_{2,i}(P)$ of
$P=(P_1,\dots,P_N)$ satisfying \eqref{formpoints}, which are uniformly bounded as $\ve\to 0$. Then, combining  \eqref{lxofglobal} and \eqref{Lxsupport} with a Taylor expansion of \eqref{approxerrfirst} about $\Gamma(y)$, we obtain
\begin{equation*}
\begin{aligned}
   &\ve^2\mu_i^2 S\left(\Psi_{0}\right)\\
   &= \kappa_i\left(\frac{3R_ih^2+R_i^3}{2h\left(h^2+R_i^2\right)^{\frac3 2}}\right)\ve\mu_i y_1\Gamma^{\gamma}_{+} +\kappa_i\gamma\Gamma^{\gamma-1}_{+}\Bigg[\nu'(1)\log\left(1+\frac{\mu^{*}_i}{\mu^{0}_i}\right)+\left(c_{1,i}\ve\mu_iy_1+c_{2,i}\ve^2\mu_i^2|y|^2\right)\Gamma(y)\\
   &+\ve\mu_i y_1|\log\ve|\widehat{\mathcal{A}}_{1,i}(P)+\ve\mu_i y_2|\log\ve|\widehat{\mathcal{A}}_{2,i}(P)+\mathcal{O}\left(\ve^2\mu_i^2|\log\ve|^2 |y|^2\right)\Bigg]+\mathcal{O}\left(\ve^2\mu_i^2|\log\ve|^2|y|^2\Gamma^{\gamma-2}_{+}\right),
    \end{aligned}
\end{equation*}
where we used \eqref{rescaledequation123}, \eqref{scalingpert}, and the expression $|x|^2=|P_i|^2+\frac{2R_i}{\sqrt{h^2+R_i^2}}\ve\mu_i y_1 +\ve^2\mu_i^2\left( \frac{h^2}{h^2+R_i^2}y_1^2 +y_2^2\right).$

One can then derive the estimate 
\begin{equation*}
\ve^2\mu_i^2 S(\Psi_{0})=\mathcal{O}\left(\ve\mu_i|\log\ve|\right).
\end{equation*}
\text{\bf Case 2:} $1 \leq |y| \leq 2.$
\\
In this annulus, we observe that $L_x(\Psi_0)=0$ due to \eqref{Lxsupport}, while \eqref{cutoffj} again gives $\eta_i=1$ and $\eta_j=0$ for every $j\neq i$. 
Therefore, taking into account  \eqref{scalingpert} and that $\Gamma(1)=0$, the error function evaluated at $|y|=1$ reduces to
\begin{equation*}
\ve^2\mu_i^2 S(\Psi_0) =\kappa_i\left(\nu'(1)\log\Big(1+\frac{\mu_i^{*}}{\mu_i^0}\Big)+\mathcal{O}(\ve\mu_i|\log\ve|)\right)^{\gamma}_{+} =\mathcal{O}\left(\ve\mu_i|\log\ve|\right).
\end{equation*}
For $1<|y|\leq 2$, we instead have $\Gamma(y)=\nu'(1)\log|y|<0,$ hence we obtain
\[\ve^2\mu_i^2S(\Psi_0)=\kappa_i\Big(\nu'(1)\log|y|\big(1+\mathcal{O}(\ve\mu_i)\big)+\mathcal{O}\left(\ve\mu_i|\log\ve|\right)\Big)^{\gamma}_{+}=\mathcal{O}\left(\ve\mu_i|\log\ve|\right).\]
\\
\textbf{Case 3}: $2 < |y| \leq \frac{\delta^2}{\ve\mu_i|\log\ve|}.$
\\
Similarly to the previous case, \eqref{cutoffj} gives $\eta_{i}=1$ and $\eta_j=0$ for all $j\neq i$, while \eqref{Lxsupport}
yields $L_x(\Psi_0)=0$. Therefore, the error function takes the form
\begin{equation*}
\begin{aligned}
\ve^2\mu_i^2 S\left(\Psi_{0}\right)&=\kappa_i\Big(\nu'(1)|\log\ve|\mathcal{K}_{1,i}(y)\Big)^{\gamma}_{+},
\end{aligned}
\end{equation*}
with
\begin{equation}\label{functK1}
\begin{aligned}
\mathcal{K}_{1,i}(y)&=\frac{\log|y|}{|\log\ve|}\left(1+c_{1,i}\ve\mu_i y_1 +c_{2,i}\ve^2\mu_i^2|y|^2\right)+\frac{\log\left(1+\frac{\mu_i^{*}}{\mu_i^0}\right)}{|\log\ve|}+\frac{1}{\nu'(1)}\left(\ve\mu_i y_1 \widehat{\mathcal{A}}_{1,i}(P)+\ve\mu_i y_2\widehat{\mathcal{A}}_{2,i}(P)\right)\\
&\hspace{3mm}+\mathcal{O}\left(\ve^2\mu_i^2|\log\ve||y|^2\right)\\
&=\frac{\log|y|}{|\log\ve|}\left(1+\mathcal{O}(\delta^2|\log\ve|^{-1})\right) + \mathcal{O}\left(\delta^2|\log\ve|^{-1}\right).
\end{aligned}
\end{equation}
In particular, \eqref{functK1} shows that there exists a small $\ve_0>0$ such that $\mathcal{K}_{1,i}(y)>0$ for all $\ve \in (0,\ve_0)$. Since $\nu'(1)|\log\ve|<0$, we deduce that  
\[\ve^2\mu_i^2S(\Psi_0)=0.\]
\text{\bf Case 4:} $\frac{\delta^2}{\ve\mu_i|\log\ve|} < |y| < \frac{\delta}{\ve\mu_i|\log\ve|}.$
\\
In this region, $\eta_{i}\in(0,1),\,\eta_j=0$ for every $j\neq i$, and $L_x(\Psi_0)=0$, while we also get the validity of the asymptotic expansion
\[\frac{\log|y|}{|\log\ve|}=1+\mathcal{O}\left(\frac{\log|\log\ve|}{|\log\ve|}\right).\]
It then follows that
\[\ve^2\mu_i^2S(\Psi_0)=\kappa_i\Big(\nu'(1)|\log\ve|\mathcal{K}_{2,i}(y)\Big)^{\gamma}_{+}\eta_i,\]
with   \[\mathcal{K}_{2,i}(y)=\frac{\log|y|}{|\log\ve|}\left(1+\mathcal{O}(\delta|\log\ve|^{-1})\right)+\mathcal{O}\left(\delta|\log\ve|^{-1}\right).\]
Using again $\nu'(1)|\log\ve|<0$, for all sufficiently small $\ve>0$ we infer that 
\[\ve^2\mu_i^2S(\Psi_0)=0.\]
\\
\text{\bf Case 5:} $\left|\frac{A_i^{-1}(x-P_i)}{\ve\mu_i}\right| \geq \frac{\delta}{\ve\mu_i|\log\ve|}, \,\, \mbox{for all} \,i=1,\dots,N.$
\\
In this unbounded region, it immediately holds 
\begin{equation*}
S(\Psi_0)=0,
\end{equation*}
since $L_{x}(\Psi_0)=0$ and $\eta_i=0$ for all $i=1,\dots,N.$
\end{proof}

\section{Inner-Outer Gluing Procedure}\label{Section 6}
In this section, we consider the approximate stream function $\Psi_{0}$ in \eqref{globalapprox} satisfying the error estimates of Proposition \ref{errorprop1}, and seek an exact solution $\Psi$ of the equation 
\begin{equation}\label{truesol}
  S\left(\Psi\right)(x)=L_{x}\left(\Psi\right)+F\left(\Psi-\frac{\alpha}{2}|\log\ve||x|^2\right)=0 \quad \text{in} \quad \R^2, 
\end{equation}
where the nonlinearity $F$ reads
\[F(s)=\sum_{j=1}^{N}\frac{\kappa_j}{\ve^2\mu_j^2}\left(\frac{s}{\kappa_j} +\nu'(1)|\log\ve|+\frac{\alpha}{2\kappa_j}|\log\ve||P_j|^2\right)^{\gamma}_{+}\eta_{j}\left(\frac{s}{\kappa_j}\right),\]
with $\gamma>3,\,u_+\defeq\max(0,u),$ and $\eta_{j},\,j=1,\dots,N$ the smooth cut-off functions described in \eqref{cutoffj}. 

We solve \eqref{truesol} by means of an \emph{Inner–Outer gluing scheme}, through which we aim to find a true solution $\Psi$ as a small perturbation of the approximate solution $\Psi_0$. We recall that $\Psi_0$ depends on the concentration points $P_1,\dots,P_N$ satisfying   \eqref{formpoints}, and on the scaling parameters $\mu_j=\mu_j^{0}+\mu_j^{*}, \, j=1,\dots,N,$ as stated in Proposition \ref{propscaling}.

More specifically,
we seek a solution $\Psi(x)$
to \eqref{truesol} of the form
\begin{equation}\label{gluingansatz}
\Psi(x)=\Psi_0(x)+\varphi(x), \quad \varphi(x)=\sum_{i=1}^{N}\tilde{\eta}_i(x)\phi_{in,i}(y_i) +\phi_{out}(x), \quad y_i=\frac{A_i^{-1}(x-P_i)}{\ve\mu_i},
\end{equation}
so that for $0<2\delta_1<\delta^2$ with $\delta$ chosen according to \eqref{cond1delta} and $\eta(x)$ given by \eqref{defeta}, the function $\tilde{\eta}_{i}(x)$ in \eqref{gluingansatz} is an additional smooth cut-off defined as
\begin{equation}\label{cutoffgluing}
\tilde{\eta}_i=\eta\left(\frac{|\log\ve|^2|A_i^{-1}(x-P_i)|}{\delta_1}\right).
\end{equation}
In view of the decomposition of the perturbation $\varphi(x)$ in \eqref{gluingansatz}, for the operator in \eqref{truesol} we write 
\begin{equation}\nonumber
\begin{aligned}
S(\Psi_0+\varphi)
 &=0 \inn \,\,\R^2,\\
 S\left(\Psi_0+\varphi\right)&=\sum_{i=1}^N\tilde{\eta}_i\bigg[L_x[\phi_{in,i}]
 +F'\left(\Psi_0-\frac{\alpha}{2}|\log\ve||x|^2\right)(\phi_{in,i}+\phi_{out})
 +S(\Psi_0)
 +N_0(\varphi)\bigg]   \\
&\hspace{4mm} + L_x[\phi_{out}]+\Big(1-\sum_{i=1}^N\tilde{\eta}_i\Big)
\bigg[F'\left(\Psi_0-\frac{\alpha}{2}|\log\ve||x|^2\right)\phi_{out}+S(\Psi_0)
 +N_0(\varphi)\bigg] \\
&\hspace{4mm} + \sum_{i=1}^N\big(L_x[\tilde{\eta}_i\phi_{in,i}]-\tilde{\eta}_iL_x[\phi_{in,i}]\big),
\end{aligned}
\end{equation}
where we introduced the notation
\begin{equation}\label{nol}
\begin{aligned}
 N_{0}(\varphi) &=F\left(\Psi_0-\frac{\alpha}{2}|\log\ve||x|^2+\varphi\right)-F\left(\Psi_0-\frac{\alpha}{2}|\log\ve||x|^2\right)-F'\left(\Psi_0-\frac{\alpha}{2}|\log\ve||x|^2\right)\varphi,\\
 F'(s)&=\sum_{i=1}^N\frac{\gamma}{\ve^2\mu_i^2}\left(\frac{s}{\kappa_i}+\nu'(1)|\log\ve|+\frac{\alpha}{2\kappa_i}|\log\ve||P_i|^2\right)^{\gamma-1}_{+}\eta_{i}\left(\frac{s}{\kappa_i}\right)\\
 &\hspace{3mm}+\sum_{i=1}^{N}\frac{1}{\ve^2\mu_i^2}\left(\frac{s}{\kappa_i}+\nu'(1)|\log\ve|+\frac{\alpha}{2\kappa_i}|\log\ve||P_i|^2\right)^{\gamma}_{+}\eta'_{i}\left(\frac{s}{\kappa_i}\right).  
\end{aligned}
\end{equation}
Rearranging terms, we deduce that $\Psi$ in \eqref{gluingansatz} is an exact solution of \eqref{truesol}, provided that the components $(\phi_{in},\phi_{out})\defeq (\phi_{in,1},\dots,\phi_{in,N},\phi_{out})$ defining $\varphi$ satisfy the coupled system of equations
\begin{align} \label{inner1}
L_x[\phi_{in,i}]+F'\left(\Psi_0-\frac{\alpha}{2}|\log\ve||x|^2\right)(\phi_{in,i}+\phi_{out})
+S(\Psi_0)+N_0\Big(\sum_{i=1}^{N}\tilde{\eta}_i\phi_{in,i}+\phi_{out}\Big)=0,&&\\
\nonumber
\mathrm{for}
\ |A_i^{-1}(x-P_i)|<\frac{2\delta_1}{|\log\ve|^2}, \quad i=1,\dots,N,&&
\end{align}
and
 \begin{equation}
     \begin{aligned} \label{outer1}
 L_x[\phi_{out}]&+\Big(1-\sum_{i=1}^N\tilde{\eta}_i\Big)\bigg[F'\left(\Psi_0-\frac{\alpha}{2}|\log\ve||x|^2\right)\phi_{out}+S(\Psi_0)
 +N_0\Big(\sum_{i=1}^{N}\tilde{\eta}_i\phi_{in,i}+\phi_{out}\Big)\bigg]\\
 &+ \sum_{i=1}^N\left(L_x[\tilde{\eta}_i\phi_{in,i}]-\tilde{\eta}_iL_x[\phi_{in,i}]\right)=0,\quad\mbox{for}\,\,x\in\mathbb{R}^2.
\end{aligned}
\end{equation}
In the sequel, we refer to \eqref{inner1} as the Inner problem and to \eqref{outer1} as the Outer problem. 

To analyse \eqref{inner1} near a vortex point 
$P_i$, we work in the rescaled variable $y=\frac{A_i^{-1}(x-P_i)}{\ve\mu_i}$, with $\mu_i=\mu_i^{0}+\mu_i^{*}$ satisfying \eqref{scalingmain} and \eqref{scalingpert}. Using Proposition \ref{expansionofL}, in the region \[|y|<\frac{2\delta_1}{\ve\mu_i|\log\ve|^2}\]  we obtain \[\ve^{2}\mu_i^2 L_{x}[\phi_{in,i}]=\Delta_{y}\phi_{in,i} +\widetilde{B}_{0}(y)[\phi_{in,i}],\] 
where $\widetilde{B}_{0}(y)[\phi_{in,i}]=\ve^2\mu_i^2 B_{0}\left(\ve \mu_i y\right)[\phi_{in,i}]$;\,see \eqref{oper}.

Moreover, recalling the definition of the functions $\widehat{\mathcal{A}}_{1,i}(P),$ and $\widehat{\mathcal{A}}_{2,i}(P)$ in \eqref{logmode1fun}, in the same region we get the expansion \[\ve^2\mu_i^2 F'\left(\Psi_{0}-\frac{\alpha}{2}|\log\ve||x|^2\right)=\gamma\Gamma^{\gamma-1}_{+}(y)+b_i(y),\] 
where
\begin{equation}\label{defofbi}
\begin{aligned}
b_i(y)&=\gamma(\gamma-1)\Gamma^{\gamma-2}_{+}\Bigg[\nu'(1)\log\left(1+\frac{\mu_i^{*}}{\mu_i^{0}}\right)+\ve \mu_i y_1\left(c_{1,i}\Gamma(y)+|\log\ve|\widehat{\mathcal{A}}_{1,i}(P)\right)
+\ve\mu_i y_2|\log\ve|\widehat{\mathcal{A}}_{2,i}(P)\\
&\hspace{28mm}+\mathcal{O}\left(\ve^2\mu_i^2|\log\ve|^2|y|^2\right)\Bigg]+\mathcal{O}\left(\ve^2\mu_i^2|\log\ve|^2|y|^2\Gamma^{\gamma-3}_{+}\right)\\
&=\mathcal{O}\left(\ve\mu_i|\log\ve|\Gamma^{\gamma-2}_{+}\right).
\end{aligned}
\end{equation}
For the nonlinear quadratic term, similar reasoning yields 
\[
\mathcal N_i(\varphi)\defeq \ve^2\mu_i^2 N_0(\varphi)
=\mathcal O\!\left(\bigl(\Gamma_+^{\gamma-2}(y)+|b_i(y)|\bigr)\varphi^2\right),
\]
while Proposition \ref{errorprop1} further implies 
\begin{equation*}
\widetilde{E}_{i}\defeq \ve^2\mu_i^2S(\Psi_0)=\mathcal{O}\left(\ve\mu_i|\log\ve|\Big(\Gamma\left(\frac{y}{2}\right)\Big)^{\gamma}_{+}\right).
\end{equation*}
Consequently, after a multiplication of \eqref{inner1} by $\ve^2\mu_i^2$, the Inner problem can be recast as
\begin{equation}\label{inprobexp}
\Delta_{y}\phi_{in,i}+\gamma\Gamma^{\gamma-1}_{+}\phi_{in,i}+\bar B_i(y)[\phi_{in,i}]+H_i(\phi_{in},\phi_{out})=0 \quad\mbox{in}\,\,B_{\rho}, \quad i=1,\dots,N,
\end{equation}
where $\rho=\frac{2\delta_1}{\ve\mu_i|\log\ve|^2}$, 
\begin{equation}\label{defbarBinner}
\bar{B}_i(y)[\phi_{in,i}]=\ve^2\mu_i^2 B_0(\ve\mu_i y)[\phi_{in,i}]+b_i(y)\phi_{in,i}(y)
\end{equation}
as in \eqref{oper} and \eqref{defofbi}, and
\begin{equation}\label{defHinner}
H_i(\phi_{in},\phi_{out})=\mathcal{N}_i\Big(\sum_{j=1}^{N}\tilde{\eta}_j\phi_{in,j}+\phi_{out}\Big)+\widetilde{E}_i+\left(\gamma\Gamma^{\gamma-1}_{+}+b_i\right)\phi_{out}.
\end{equation}
\medskip
Regarding the Outer problem in \eqref{outer1}, for notational simplicity we write
\begin{equation}\label{Poissonouter}
L_{x}\left[\phi_{out}\right]+G\left(\phi_{in},\phi_{out}\right)= 0 \quad\mbox{in} \quad\R^2,
\end{equation}
where for $\phi_{in}=(\phi_{in,1},\dots,\phi_{in,N})$ we have introduced 
\begin{equation}\label{Gout}
\begin{aligned}
G\left(\phi_{in},\phi_{out}\right)=&\Big(1-\sum_{i=1}^{N}\tilde{\eta}_i\Big)\bigg[F'\left(\Psi_{0}-\frac{\alpha}{2}|\log\ve||x|^2\right)\phi_{out}+S\left(\Psi_{0}\right)+N_0\Big(\sum_{i=1}^{N}\tilde{\eta}_i\phi_{in,i}+\phi_{out}\Big)\bigg]\\
&+\sum_{i=1}^N\big(L_x[\tilde{\eta}_i\phi_{in,i}]-\tilde{\eta}_iL_x[\phi_{in,i}]\big).
\end{aligned}
\end{equation}
\subsection{Solution of the Outer Problem}\label{subsecouter}
To initiate the resolution of the coupled system formed by \eqref{inprobexp} and \eqref{Poissonouter}, we first deal with the Outer problem \eqref{Poissonouter}. For $\sigma>0$ and $\beta\in(0,1)$, we make the a priori assumptions that 
\begin{equation}\label{b100}
\phi_{out}\in C^{1,\beta}(\R^2), \quad\|(1+|x|^2)^{-1
}\phi_{out}\|_{\infty}\leq 100,
\end{equation}
while $\phi_{in,i}$ satisfies
\begin{equation}\label{propertiesforphi1}
(1+|y|)|D_y \phi_{in,i}(y)|+|\phi_{in,i}(y)| \leq \frac{C\ve\mu_i|\log\ve|}{1+|y|^{\sigma}},\quad i=1,\dots,N,
\end{equation}
with $\mu_i$ as in Proposition \ref{propscaling} and  $y=\frac{A_i^{-1}(x-P_i)}{\ve\mu_i}$.

In this setting, the following lemma holds.
\begin{lemma}\label{lemmaout}
Let $\phi_{out}(x)$ as in  \eqref{b100} and $\phi_{in,i}(y),\,i=1,\dots,N,$ satisfy the estimate \eqref{propertiesforphi1}. Then, for $\delta>0$ sufficiently small and the cut-off functions $\eta_{i}, \,\tilde{\eta}_{i},\,i=1,\dots,N,$ in \eqref{cutoffj} and \eqref{cutoffgluing}, there exists a small $\ve_0 >0$ such that for all $\ve \in (0,\ve_0)$ it holds \[G\left(\phi_{in},\phi_{out}\right)=\sum_{i=1}^{N}\big(L_{x}\left[\tilde{\eta}_i\phi_{in,i}\right]-\tilde{\eta}_iL_{x}\left[\phi_{in,i}\right]\big),\]
where $G(\phi_{in},\phi_{out})$  is defined in \eqref{Gout}.
\end{lemma}
\begin{proof}
For fixed $i\in\{1,\dots,N\}$, we first consider the inner region \[|A_i^{-1}(x-P_i)|\leq\frac{\delta_1}{2|\log\ve|^2}.\] 
Due to \eqref{cutoffj}, \eqref{cutoffgluing} and the uniform separation of the points in \eqref{unifdistance}, namely \[|P_j-P_i| \geq \frac{d}{|\log\ve|}, \quad \mbox{with} \quad 0<2\delta_1<\delta^2 <\delta < \frac{\sqrt{h^2+r_0^2}}{4h}d,\] 
we have $\tilde{\eta}_i=1$ and $\tilde{\eta}_j=0$ for every $j\neq i$, hence the result follows directly since $1-\sum\limits_{j=1}^{N}\tilde{\eta}_j=1-\tilde{\eta}_i= 0.$

To establish the lemma, it therefore suffices to examine \eqref{Gout} in the region \[|A_i^{-1}(x-P_i)|>\frac{\delta_1}{2|\log\ve|^2}.\]  
We recall that in the definition of the nonlinearity $F$ in \eqref{nonlinearity}, the cut-off $\eta_i$ was introduced in \eqref{cutoffj}-\eqref{rangecutoffi} through $\frac{1}{\kappa_i}\left(\Psi_0-\frac{\alpha}{2}|\log\ve|\,|x|^2\right)=\mathcal{O}(|\log\ve|),$ and the transition occurs on a $\mathcal{O}(1)$ scale. However, the same cut-off $\eta_i$ is now evaluated at the perturbed argument $\frac{1}{\kappa_i}(\Psi_0-\frac{\alpha}{2}|\log\ve||x|^2+\tilde{\eta}_i\phi_{in,i}+\phi_{out})$. By \eqref{b100} and \eqref{propertiesforphi1}, the perturbation satisfies
$\frac{1}{\kappa_i}(\tilde{\eta}_i\phi_{in,i}+\phi_{out})=\mathcal{O}(1)$ in any bounded region, whereas on unbounded sets, the term $-\frac{\alpha}{2}|\log\ve||x|^2<0$ dominates the quadratic growth $\mathcal{O}(1+|x|^2)$ of the global correction $H_{2\ve}(x)$ appearing in $\Psi_0$ and the outer perturbation $\phi_{out}(x)$ (see \eqref{globalapprox} and \eqref{boundh2e}) for all small $\ve>0.$ Consequently, the transition region of $\eta_i$ differs from that in \eqref{cutoffj} only up to multiplicative constants, i.e. there exist universal constants $\lambda_1,\lambda_2>0$ (independent of $\ve>0$) such that
\begin{equation*}
\eta_i\left(
\frac{1}{\kappa_i}\left(
\Psi_0-\frac{\alpha}{2}|\log \ve|\,|x|^2+\tilde{\eta}_i\phi_{in,i}+\phi_{out}
\right)\right)
=
\begin{cases}
1, & \text{if } \left|A_i^{-1}(x-P_i)\right| \le \dfrac{\lambda_1\delta^2}{|\log \ve|},\\[1ex]
0, & \text{if } \left|A_i^{-1}(x-P_i)\right| \ge \dfrac{\lambda_2\delta}{|\log \ve|}.
\end{cases}
\end{equation*}
In this respect, for the analysis of \eqref{Gout} we focus on the regions
\begin{equation}\label{Hregions}
\begin{aligned}
U_{in,i}\defeq\left\{\frac{\delta_1}{2|\log\ve|^2} <|A_{i}^{-1}(x-P_i)|<\frac{\lambda_2\delta}{|\log\ve|}\right\},\quad
U_{out}\defeq \bigcap\limits_{j=1}^{N}\left\{|A_j^{-1}(x-P_j)|\geq\frac{\lambda_2\delta}{|\log\ve|}\right\},
\end{aligned}
\end{equation}
keeping in mind that for $x\in U_{in,i},$ one has $\eta_j=\tilde{\eta}_j=0$ for every $j\neq i$. 

Let us first consider $x\in U_{in,i}.$ By Proposition \ref{errorprop1} and \eqref{nol}, we find \[\left(1-\tilde{\eta}_i\right)\Big(F'\big(\Psi_0-\frac{\alpha}{2}|\log\ve||x|^2\big)\phi_{out}+S(\Psi_0)\Big)=0,\]
since 
\begin{equation}\label{errorandfsupport}
\operatorname{supp} F'\left(\Psi_0-\frac{\alpha}{2}|\log\ve||x|^2\right)\,\cup\, \operatorname{supp} S(\Psi_0) \subset \bigcup\limits_{j=1}^{N}\left\{x\in\R^2: |A_j^{-1}(x-P_j)| < 2\ve\mu_j\right\}.
\end{equation}
Therefore, the only remaining term to be controlled in the region $U_{in,i}$ is 
\begin{equation*}
\begin{aligned}
&(1-\tilde{\eta}_i)N_0(\varphi)=(1-\tilde{\eta}_i)F\Big(\Psi_0-\frac{\alpha}{2}|\log\ve||x|^2+\tilde{\eta}_i\phi_{in,i}+\phi_{out}\Big)\\
&=\frac{(1-\tilde{\eta}_i)\kappa_i}{\ve^2\mu_i^2}\left(\frac{1}{\kappa_i}\left(\Psi_0-\frac{\alpha}{2}|\log\ve||x|^2+\tilde{\eta}_i\phi_{in,i}+\phi_{out}\right)+\nu'(1)|\log\ve|+\frac{\alpha}{2\kappa_i}|\log\ve||P_i|^2\right)^{\gamma}_{+}\\
&\hspace{6mm}\times\eta_i\left(\frac{1}{\kappa_i}\left(\Psi_0-\frac{\alpha}{2}|\log\ve||x|^2+\tilde{\eta}_i\phi_{in,i}+\phi_{out}\right)\right).
\end{aligned}
\end{equation*} 
For readability, we set
\[\mathcal{H}_i(x)\defeq(1-\tilde{\eta}_i)F\Big(\Psi_0-\frac{\alpha}{2}|\log\ve||x|^2+\tilde{\eta}_i\phi_{in,i}+\phi_{out}\Big),\]
and establish that $\mathcal{H}_i(x)\equiv 0$  for all small $\ve>0$ due to \eqref{b100} and \eqref{propertiesforphi1}. We split the analysis into three annular regions around $P_i$ using the variable 
 \[z=A_{i}^{-1}(x-P_i), \quad z=\ve\mu_i y,\]
 corresponding to the transition zones of the cut-offs $\eta_i$ and $\tilde{\eta}_i$.
\\
\text{\bf Case 1}:  $\frac{\delta_1}{2\ve\mu_i|\log\ve|^2}< |y| \leq \frac{\delta_1}{\ve\mu_i|\log\ve|^2}$.
\\
In this region, $\tilde{\eta}_i\in(0,1)$ and  $\eta_{i}=1$, thus Proposition \ref{errorprop1} and the definitions of $\widehat{\mathcal{A}}_{1,i}(P)$ and $\widehat{\mathcal{A}}_{2,i}(P)$ in \eqref{mode1func}-\eqref{logmode1fun} allow us to write 
\begin{equation*}
\begin{aligned}
\mathcal{H}_i(x)
&=\frac{\left(1-\tilde{\eta}_i\right)\kappa_i}{\ve^2\mu_i^2}\Bigg(\Gamma(y)(1+c_{1,i}\ve\mu_iy_1+c_{2,i}\ve^2\mu_i^2|y|^2)+\nu'(1)\log\left(1+\frac{\mu_i^{*}}{\mu_i^{0}}\right)+\ve\mu_iy_1|\log\ve|\widehat{\mathcal{A}}_{1,i}(P)\\
&\hspace{4mm}+\ve\mu_i y_2 |\log\ve| \widehat{\mathcal{A}}_{2,i}(P)
+\mathcal{O}\left(\ve^2\mu_i^2|\log\ve|^2|y|^2\right)+\frac{1}{\kappa_i}\big(\tilde{\eta}_i\phi_{in,i}(y)+\phi_{out}\left(P_i+A_i\ve\mu_i y\right)\big)\Bigg)^{\gamma}_{+}\\
&=\frac{\left(1-\tilde{\eta}_i\right)\kappa_i}{\ve^2\mu_i^2}\Big(\nu'(1)|\log\ve|\mathcal{E}_{1,i}(y)\Big)^{\gamma}_{+},
\end{aligned}
\end{equation*}
with
\begin{equation*}
\begin{aligned}
\mathcal{E}_{1,i}(y)&=\frac{\log|y|}{|\log\ve|}\left(1+c_{1,i}\ve\mu_iy_1+c_{2,i}\ve^2\mu_i^2|y|^2\right)+\frac{\log\left(1+\frac{\mu_i^{*}}{\mu_i^0}\right)}{|\log\ve|}+\frac{\ve\mu_i}{\nu'(1)}\Big(y_1\widehat{\mathcal{A}}_{1,i}(P)+y_2\widehat{\mathcal{A}}_{2,i}(P)\Big)\\
&\hspace{4mm}+\mathcal{O}\left(\ve^2\mu_i^2|\log\ve||y|^2\right)+\frac{1}{\kappa_i\nu'(1)|\log\ve|}\Big(\tilde{\eta}_{i}\phi_{in,i}(y)+\phi_{out}\left(P_i+A_i\ve\mu_i y\right)\Big).
\end{aligned}
\end{equation*}
Using \eqref{b100} and \eqref{propertiesforphi1} together with the expansion $\frac{\log|y|}{|\log\ve|}= 1+\mathcal{O}\left(\frac{\log|\log\ve|}{|\log\ve|}\right),$ we obtain \[\mathcal{E}_{1,i}(y)=\left(1+\mathcal{O}\left(\frac{\log|\log\ve|}{|\log\ve|}\right)\right)\left(1+\mathcal{O}\left(\delta_1|\log\ve|^{-2}\right)\right)+\mathcal{O}\left(|\log\ve|^{-1}\right),\] while recalling that $\nu'(1)|\log\ve|<0$, we may take  $\ve_0>0$ sufficiently small to ensure $\mathcal{E}_{1,i}(y)>0$, so that $\mathcal{H}_i(x)=0$ for all $\ve\in(0,\ve_0)$. 
\\
\text{\bf Case 2}: $\frac{\delta_1}{\ve\mu_i|\log\ve|^2} < |y| \leq \frac{\lambda_1\delta^2}{\ve\mu_i|\log\ve|}$. 
\\
In this region, $\tilde{\eta}_i=0$ and $\eta_i=1$, hence we get 
\begin{equation*}
\begin{aligned}
\mathcal{H}_i(x)&=\frac{\kappa_i}{\ve^2\mu_i^2}\Bigg(\Gamma(y)(1+c_{1,i}\ve\mu_i y_1 +c_{2,i}\ve^2\mu_i^2|y|^2)+\nu'(1)\log\left(1+\frac{\mu_i^{*}}{\mu_i^{0}}\right)+\ve\mu_iy_1|\log\ve|\widehat{\mathcal{A}}_{i,1}(P)\\
&\hspace{15mm}+\ve\mu_i y_2 |\log\ve| \widehat{\mathcal{A}}_{i,2}(P)+\mathcal{O}\left(\ve^2\mu_i^2|\log\ve|^2|y|^2\right)+\frac{1}{\kappa_i}\phi_{out}\left(P_i+A_i\ve\mu_i y\right)\Bigg)^{\gamma}_{+}
\end{aligned}
\end{equation*}
 Analogously to the previous case, for all sufficiently small $\ve>0$ it holds $\mathcal{H}_i(x)=0$,
 since \[\mathcal{H}_i(x)=\frac{\kappa_i}{\ve^2\mu_i^2}\left(\nu'(1)|\log\ve|\mathcal{E}_{2,i}(y)\right)^{\gamma}_{+},\]
 with
 \[\mathcal{E}_{2,i}(y)=
 \left(1+\mathcal{O}\Big(\frac{\log|\log\ve|}{|\log\ve|}\Big)\right)\left(1+\mathcal{O}\left(\delta^2|\log\ve|^{-1}\right)\right) +\mathcal{O}\left(|\log\ve|^{-1}\right).\]
\\
\text{\bf Case 3:} $\frac{\lambda_1\delta^2}{\ve\mu_i|\log\ve|} <|y| < \frac{\lambda_2\delta}{\ve\mu_i|\log\ve|}.$
\\
Here $\eta_{i}\in(0,1)$ and $\tilde{\eta}_i=0$, and a similar argument as in the previous two cases establishes $\mathcal{H}_i(x)=0$ for all $\ve>0$ sufficiently small.


We now turn our attention to the outer region $U_{out}$ in \eqref{Hregions}. By definition, $\tilde{\eta}_i= 0$ and $\eta_i=0$ for all $i\in\{1,\dots,N\}$, hence using the support property \eqref{errorandfsupport} we directly have
\begin{equation*}
\bigg(1-\sum_{i=1}^{N}\tilde{\eta}_i\bigg)\left[F'\left(\Psi_{0}-\frac{\alpha}{2}|\log\ve||x|^2\right)\phi_{out}+S\left(\Psi_{0}\right)+N_0\Big(\sum_{i=1}^{N}\tilde{\eta}_i\phi_{in,i}+\phi_{out}\Big)\right]=0.
\end{equation*}
\end{proof}
\begin{corollary}
A key outcome of Lemma \ref{lemmaout} is that if $\phi_{out}$ and $\phi_{in,i},\,i=1,\dots,N,$ satisfy \eqref{b100} and \eqref{propertiesforphi1}, then the Outer problem \eqref{Poissonouter} reduces to 
\begin{equation}\label{poisreduced}
L_{x}[\phi_{out}]=-\sum_{i=1}^{N}\big(L_{x}\left[\tilde{\eta}_i\phi_{in,i}\right]-\tilde{\eta}_i L_{x}\left[\phi_{in,i}\right]\big) \quad \textnormal{in} \quad \R^2.
\end{equation}
\end{corollary}
As a consequence, provided that the right-hand side of \eqref{poisreduced} decays sufficiently fast at infinity, the problem falls within the class of equations treated in Proposition \ref{propouter}, and therefore the existence of a corresponding solution $\phi_{out}(x)$ follows directly.

To verify the decay property, in the variable \[z_i=A_i^{-1}(x-P_i), \, z_i=\ve\mu_i y_i,\]
a direct substitution in \eqref{oper} gives 
\begin{equation}\label{outeroperatorphiin}
\begin{aligned}
&\sum_{i=1}^{N}\big(L_{x}\left[\tilde{\eta}_i\phi_{in,i}\right]-\tilde{\eta}_i L_{x}\left[\phi_{in,i}\right]\big)\\
&=\sum_{i=1}^{N}\Bigg[2\nabla_{z_i}\tilde{\eta}_{i}\left(\frac{|\log\ve|^2|z_i|}{\delta_1}\right)\cdot\nabla_{z_i}\phi_{in,i}(y_i)+\phi_{in,i}(y_i)\Delta_{z_i}\tilde{\eta}_i\left(\frac{|\log\ve|^2|z_i|}{\delta_1}\right)\\
&\hspace{4mm}+B_0(z_i)\left[\tilde{\eta}_i\left(\frac{|\log\ve|^2|z_i|}{\delta_1}\right)\phi_{in,i}(y_i)\right]-\tilde{\eta}_i\left(\frac{|\log\ve|^2|z_i|}{\delta_1}\right)B_0(z_i)\left[\phi_{in,i}(y_i)\right]\Bigg].        
\end{aligned}
\end{equation}
For each $i\in\{1,\dots,N\},$ one can deduce that the function $L_{x}[\tilde{\eta}_i\phi_{in,i}]-\tilde{\eta}_iL_{x}[\phi_{in,i}]$
is bounded and compactly supported in the transition region of the cut-off $\tilde{\eta}_i$, thus there exist $C>0$ and $\bar\nu >2$ (independent of $i$) such that \[\big|L_{x}[\tilde{\eta}_i\phi_{in,i}]-\tilde{\eta}_iL_{x}[\phi_{in,i}]\big| \leq \frac{C}{1+|x|^{\bar\nu}},\quad \forall x\in\R^2.\] 
In particular, \[\sup\limits_{x\in\R^2}\big|L_{x}\left[\tilde{\eta}_i\phi_{in,i}\right]-\tilde{\eta}_iL_{x}\left[\phi_{in,i}\right]\big|\leq C\ve^{1+\sigma^{*}}\]
uniformly in $i=1,\dots,N$, for some $0<\sigma^{*}<\sigma$, with $\sigma>0$ as in \eqref{propertiesforphi1}. 

To illustrate the estimate above, using \eqref{propertiesforphi1}, \eqref{cutoffgluing} and the variable $z=A_i^{-1}(x-P_i),\,z=\ve\mu_i y$, we have
\begin{equation}\nonumber
\begin{aligned}
\Bigg|\nabla_{z}\tilde{\eta}_i\left(\frac{|\log\ve|^2|z|}{\delta_1}\right)\cdot\nabla_{z}\phi_{in,i}(y)\Bigg|&\leq \frac{C|\log\ve|^2
}{\ve\mu_i}\tilde{\eta}_i^{'}\left(\frac{|\log\ve|^2|z
|}{\delta_1}\right)\nabla_{y}\phi_{in,i}(y)\\
&\leq o(1) \frac{\ve ^{1+\sigma^{*}}}{1+|x|^{\bar{\nu}}},
\end{aligned}
\end{equation}
with $o(1)\to 0$ as $\ve \to 0,\, \bar{\nu}>2$ and some $0<\sigma^{*}<\sigma$, while the remaining terms in \eqref{outeroperatorphiin} can be handled in a similar manner.

As a result, we apply Proposition \ref{propouter} to get the existence of a solution $\phi_{out}\in C^{1,\beta}(\R^2)$ for any $\beta\in(0,1)$ to equation \eqref{poisreduced}, which satisfies 
\begin{equation}\label{boundphiout}
|\phi_{out}(x)|\leq C \ve^{1+\sigma^{*}}\left(1+|x|^2\right),\quad \forall x \in \R^2.
\end{equation}
Finally, since $\phi_{out}(x)$ is unique up to a constant, we choose the one with the property
\begin{equation}\label{phioutP}
\phi_{out}(r_0,0)=0.
\end{equation}
\section{Inner and Outer Linear Theories}\label{Section 7}
This section collects the linear solvability results and a priori estimates for the Inner and Outer problems introduced in Section \ref{Section 6}. These constitute the main analytic tools for solving the coupled system in \eqref{inprobexp} and \eqref{Poissonouter}. Although the linear theory for the Inner problem was established in \cite{MR5036920}, we include a brief sketch of the proofs in the Appendix for the reader's convenience.
\subsection{Linear Inner Theory}
Let $h(y)$ be a function satisfying the decay condition 
\begin{equation}\label{hdecay}
    |h(y)|\leq \frac{C}{1+|y|^{2+\sigma}},
\end{equation} for some $\sigma>0$. The current section considers the linear equation
\be\label{Linearisedinner}
\begin{aligned}
\Delta_{y} \phi + \gamma\Gamma^{\gamma-1}_{+}\phi +  h(y)= 0 \inn \, \R^2, 
\end{aligned}
\ee
where $\Gamma$ is given in \eqref{defGamma}.

For $\sigma>0$ as in \eqref{hdecay} and $\beta\in(0,1)$, we employ the weighted norms
\be\label{norma} \begin{aligned}
\|h\|_{2+\sigma}  = & \sup_{y\in \R^2} (1+|y|)^{2+\sigma}|h(y)|, \\
\|h\|_{2+\sigma,\beta}  = &\|h\|_{2+\sigma}  +   \sup_{y\in \R^2}(1+|y|)^{2+\sigma + \beta}[h]_{B_1(y),\beta} ,
\end{aligned}
\ee
where for any $A\subset \R^2,$ we use the standard H\"{o}lder seminorm notation
$$
[h]_{A,\beta} = \sup_{\ell_1\neq \ell_2, \ell_1,\ell_2\in A}  \frac {|h(\ell_1) -h(\ell_2)| } {|\ell_1-\ell_2|^\beta}.
$$
In \cite{MR2456896}, the authors established that the linear operator $\Delta_y +\gamma\Gamma^{\gamma-1}_{+}$ in \eqref{Linearisedinner} is nondegenerate in $L^{\infty}(\R^2)$, in the sense that its kernel is generated by the bounded translation modes $Z_i(y)=\partial_{y_i}\Gamma(y), \ i=1,2$. On the other hand, if one enlarges the class of admissible solutions of \eqref{Linearisedinner} to allow logarithmic growth at infinity, then the kernel becomes three-dimensional, and it is spanned by
\begin{equation}\label{elemkernel}
Z_0(y)=\frac{2}{\gamma-1}\Gamma(y)+y\cdot\nabla_y \Gamma(y), \quad  Z_i(y)  =  \partial_{y_i}  \Gamma(y), \  i=1,2.
\end{equation}
Further details on the radial scaling mode $Z_0$ can be found in \cite{MR5036920}*{Section 5.2}. Here, we note that $Z_0(y)=\mathcal{O}\left(\log\left(2+|y|\right)\right)$ and $Z_i(y)=\mathcal{O}\left(\frac{1}{1+|y|}\right),\,i=1,2$. 

The following lemma holds.
\begin{lemma}\label{lemat}
Given $\sigma>0$ and $\beta\in (0,1)$, consider the norms defined in \eqref{norma} and the functions $Z_i,\,i=0,1,2,$ given by \eqref{elemkernel}. For each $h$ with $\|h\|_{2+\sigma} <+\infty$,
there exist a constant $C>0$ and a solution $\phi =  \mathcal{T} [ h]$ of equation \eqref{Linearisedinner}, that defines a linear operator of $h$ and satisfies
\begin{equation*}
\begin{aligned}
&(1+|y|) |\nabla \phi (y)|  + | \phi (y)| \\
&\le  \,  C \left [ \,  \log \left(2+|y|\right) \left|\int_{\R^2} h Z_0\right|  +    (1+|y|) \sum_{i=1}^2 \left|\int_{\R^2} h Z_i\right|+  (1+|y|)^{-\sigma} \|h\|_{2+\sigma}   \,\right].
\end{aligned}
\end{equation*}
Moreover, if 
$\|h\|_{2+\sigma,\beta} <+\infty$, it further holds
\begin{equation*}
\begin{aligned}
& (1+|y|^{2+\beta})  [D^2_y \phi]_{B_1(y),\beta}  +\left(1+|y|^2\right)  \left|D^2_y \phi (y)\right| \\  &
\,  \le  \,   C \left [ \,  \log \left(2+|y|\right) \left|\int_{\R^2} h Z_0\right|  +    (1+|y|) \sum_{i=1}^2 \left|\int_{\R^2} h Z_i\right|  +  (1+|y|)^{-\sigma} \|h\|_{2+\sigma,\beta}   \,\right]. 
\end{aligned}
\end{equation*}
\end{lemma}

\subsection{Projected Linearised Inner Problem}
In this section, we fix $\delta_{1}>0$ sufficiently small, take scaling parameters $\mu_i,\,i=1,\dots,N$ as in Proposition \ref{propscaling}, and define 
\[\rho \defeq \frac{2\delta_1}{\ve\mu_i|\log\ve|^2}.\]
For certain scalars $d_{ij},\,j=0,1,2,$ and a fixed $i\in\{1,\dots,N\},$ we  consider the projected linearised problem
\begin{equation}\label{problemlinearprojected}
    \Delta_{y}\phi_{i}+\gamma\Gamma^{\gamma-1}_{+}\phi_{i}+\bar{B}_i[\phi_{i}]+h(y)=\sum_{j=0}^{2}d_{ij}\gamma\Gamma^{\gamma-1}_{+}Z_j\quad\mbox{in}\quad B_{\rho}.
\end{equation}
For convenience, we recall that the definition of the functions $Z_j, j=0,1,2$ can be found in \eqref{elemkernel}, while 
 \begin{equation}\label{defofBinner}
 \bar{B}_i[\phi_{i}]=b_i(y)\phi_{i}+\widetilde{B}_0[\phi_{i}],
 \end{equation}
 with $b_i(y)=\mathcal{O}\left(\ve\mu_i|\log\ve|\Gamma^{\gamma-2}_{+}\right)$ and $\widetilde{B}_{0}=\ve^2\mu_i^2B_{0}(\ve\mu_i y)$; see \eqref{oper} and \eqref{defofbi}. 

Moreover, for $\sigma>0$,\, $\beta\in(0,1)$, and any open set $A \subset \R^2$, we modify the norms in \eqref{norma} taking all suprema over $A$, namely
\begin{equation}\label{norm1}
\|h\|_{2+\sigma,A}=\sup_{y\in A}\left(1+|y|\right)^{2+\sigma}|h(y)|, \quad \|h\|_{2+\sigma,\beta,A}=\sup_{y\in A}\left(1+|y|\right)^{2+\sigma+\beta}[h]_{B_{1}(y)\cap A,\beta}+\|h\|_{2+\sigma,A},\end{equation}
while for $\phi_{i}\in C^{2,\beta}(A)$ we define the norm
 \begin{equation}\label{starnorm}
 \|\phi_{i}\|_{*,\sigma,A} =  \|  D^2_y\phi_{i} \|_{2+\sigma,\beta,A}   + \|  D_y\phi_{i} \|_{1+\sigma,A}+ \|\phi_{i} \|_{\sigma,A}.
 \end{equation}
With a slight abuse of notation, we will suppress the subscript in \eqref{norm1} and \eqref{starnorm} when $A=\R^2$. 

We establish the following result.
\begin{prop}\label{prop1} Let $\sigma>0$,\,$\beta \in (0,1)$, and consider the norms defined in \eqref{norm1} and \eqref{starnorm}. Then, for $\ve>0$ sufficiently small and for every $h$ satisfying $\|h\|_{2+\sigma,\beta,B_{\rho}} < +\infty$, there exists a constant $C>0$ such that for any differential operator $\bar{B}_i$ of the form \eqref{defofBinner} and $\rho=\frac{2\delta_1}{\ve\mu_i|\log\ve|^2}$, the projected
problem \eqref{problemlinearprojected} has a solution $\phi_{i} = T_i[h]$ for certain scalars $d_{ij}[h], \,j=0,1,2$, which defines a linear operator of $h$ and satisfies
\begin{align*} 
\| \phi_{i}\|_{*,\sigma, B_\rho} \ \le\ C \|h \|_{2+\sigma,\beta, B_\rho} .
\end{align*}
Moreover, the coefficients $d_{ij}$ admit the expansions
\begin{align*}
d_{i0}[h]\, = & \, \gamma_0\int_{B_{\rho}}  h Z_0  + \mathcal{O}\left(\frac{\log(2+\rho)}{1+\rho^{\sigma}}\right)\|h \|_{2+\sigma,\beta, B_\rho} , \\    d_{ij}[h]\, = &\, \gamma_j\int_{B_{\rho}}  h Z_j  + \mathcal{O}\left(\ve\mu_i|\log\ve|\right)  \|h \|_{2+\sigma,\beta, B_\rho}, \ j=1,2,
\end{align*}
where
$\gamma_j^{-1} = \int_{\R^2} \gamma\Gamma^{\gamma-1}_{+} Z_j^2,\,j=0,1,2.$ \end{prop}
\subsection{Linear Outer Theory}
We consider the differential operator $L_x$ defined in \eqref{operL}, and study the elliptic equation  
	\begin{equation}
	\label{outerprob}
	L_x[\psi] + g(x) = 0  \inn \R^2 ,
	\end{equation}
	where $g(x)$ is a bounded function satisfying the decay condition
	$$
    \|g\|_{\bar{\nu}}\, :=\, \sup_{x\in \R^2}\left(1+|x|^{\bar{\nu}}\right)|g(x)|< +\infty ,
    $$
	for some $\bar{\nu} >2$. 
    
 The following Proposition holds.
\begin{prop}\label{propouter}
		There exists a solution $\psi(x)$ to \eqref{outerprob} of class $C^{1,\beta}(\R^2)$ for any $\beta\in(0,1)$, which defines a linear operator $\psi= {\mathcal T}^o (g) $ of $g(x)$
		and satisfies the bound
		\begin{equation*}
		|\psi(x)| \,\le \,C\| g\|_{\bar{\nu}} (1+ |x|^2) \quad \forall x \in \R^2,
		\end{equation*}
		for a constant $C>0$.
	\end{prop}
\begin{proof}
See \cite{MR4417384}*{Proposition 7.1}.
\end{proof}

\section{Solvability of the Inner-Outer System}\label{Section projected}
With $\phi_{out}$ already obtained in Section \ref{subsecouter} as a solution of the Outer problem \eqref{poisreduced}, we now aim to utilise the linear solvability theories developed in Section \ref{Section 7} to construct, for each $i=1,\dots,N$, a solution $\phi_{in,i}$ for the Inner problem
\[\Delta_{y}\phi_{in,i}+\gamma\Gamma^{\gamma-1}_{+}\phi_{in,i}+\bar{B}_i[\phi_{in,i}]
+H_i\left(\phi_{in},\phi_{out}\right) = 0 \quad \mbox{in} \,\, B_{\rho}.\]
We recall the notation $\phi_{in}=\left(\phi_{in,1},\dots,\phi_{in,N}\right)$,\, and that the solution $\phi_{out}$ satisfies \eqref{boundphiout} and \eqref{phioutP}. In addition,  $\rho=\frac{2\delta_1}{\ve\mu_i|\log\ve|^2}$, while the definitions of $\Bar{B}_i$ and $H_i(\phi_{in},\phi_{out})$ appear in \eqref{defofBinner} and \eqref{defHinner}.

As a starting point, we decompose \[\phi_{in,i}=\phi_{i,1}+\phi_{i,2},\quad i=1,\dots,N,\] so that for the functions $Z_j$ in \eqref{elemkernel} and scalars $d_{ij}, j=0,1,2$ as in Proposition \ref{prop1}, the function $\phi_{i,1}$ solves the auxiliary projected problem
\begin{equation}\label{problemforphi1}
\Delta_{y}\phi_{i,1}+\gamma\Gamma^{\gamma-1}_{+}\phi_{i,1}+\bar{B}_i[\phi_{i,1}]+\bar{B}_i[\phi_{i,2}]+H_i(\phi_{in},\phi_{out})=\sum_{j=0}^{2}d_{ij}\gamma\Gamma^{\gamma-1}_{+}Z_j \quad \mbox{in} \quad B_{\rho},
\end{equation}
whereas $\phi_{i,2}$ is required to satisfy the linear equation 
\begin{equation}\label{problemforphi2}
\Delta_{y}\phi_{i,2}+\gamma\Gamma^{\gamma-1}_{+}\phi_{i,2} +d_{i0}\gamma\Gamma^{\gamma-1}_{+}Z_0=0 \quad \mbox{in}\quad\R^2.
\end{equation}
In this setting, we recognise that the Inner–Outer system in \eqref{inprobexp} and \eqref{poisreduced} admits a solution of the form \[(\phi_{in},\phi_{out}), \quad \phi_{in}=\left(\phi_{1,1}+\phi_{1,2},\dots,\phi_{N,1}+\phi_{N,2}\right),\] as long as we further attain the coupled conditions 
\begin{equation}\label{redconds}d_{ij}\left[H_i(\phi_{in},\phi_{out})+\bar{B}_i[\phi_{i,2}]\right]=0, \quad i=1,\dots,N, \, j=1,2.
\end{equation}
We postpone the analysis of \eqref{redconds} to Section \ref{reduced}, where we appropriately adjust the location of the vortex points $P_1,\dots,P_N$ in \eqref{formpoints} to achieve their solvability.

We currently examine the equation for $\phi_{i,2}$ in \eqref{problemforphi2}. Setting $y=re^{i\theta}, r=|y|,$ and denoting by $\xi_0$ the unique root of $Z_0$ in $(0,1),$ a smooth radial solution of this equation is given by
\begin{equation}\label{explicitphi2}
\begin{aligned}
\phi_{i,2}(r)&=-d_{i0} Z_0 (r) \int_{\xi_0}^{r}\frac{\dd s}{sZ_0(s)^2} \int_{0}^{s} \gamma\Gamma^{\gamma-1}_{+}(\rho) Z_{0}^2(\rho)\rho \, \dd \rho\\
&\eqcolon d_{i0}\hat{\phi}_{2}(r).
\end{aligned}
\end{equation}
Since the term
$\Gamma^{\gamma-1}_{+}$ appearing in \eqref{explicitphi2} is supported in $B_1(0)$, for any $s>1$ it holds \[\int_{0}^{s}\gamma\Gamma^{\gamma-1}_{+}(\rho)Z_0(\rho)^2\rho \, \dd \rho = \int_{0}^{1}\gamma\Gamma^{\gamma-1}_{+}(\rho)Z_0(\rho)^2\rho \, \dd \rho, \]
which in turn yields the logarithmic growth bounds 
\begin{equation}\label{estimateforphi2}
|\hat{\phi}_{2}(r)|\leq C \log(2+r), \quad |\phi_{i,2}(r)| \leq C|d_{i0}|\log(2+r), \quad \mbox{for all}\,\,r>0.
\end{equation}
 The Proposition below is the main result of this section.
\begin{prop}\label{propsec8}
Let $\sigma > 0$ and $\beta \in (0,1)$. Then, for all $\ve>0$ sufficiently small, there exist constants $C>0$ and $\sigma^{*}\in(0,\sigma$), functions $\phi_{out},\,\mathbf{\Phi}_1=\left(\phi_{1,1},\dots,\phi_{N,1}\right),\mathbf{\Phi}_2=\left(\phi_{1,2},\dots,\phi_{N,2}\right)$ and coefficients $d_{i0},\,i=1,\dots,N,$ solving \eqref{poisreduced}, \eqref{problemforphi1} and \eqref{problemforphi2}, satisfying
 \begin{equation}\label{boundsforsoln}
    \left\|(1+|x|^2)^{-1}\phi_{out}\right\|_{\infty}\leq C\ve^{1+\sigma^{*}}, \quad \|\phi_{i,1}\|_{*,\sigma,B_{\rho}}\leq C\ve\mu_i|\log\ve|, \quad |d_{i0}|\leq C (\ve\mu_i)^{1+\sigma}|\log\ve|^{2+2\sigma},
 \end{equation}
 for all $i=1,\dots,N$.
 
 In the above, $\mu_i$ denote the scaling parameters of Proposition \ref{propscaling}, \,$\rho=\frac{2\delta_1}{\ve\mu_i|\log\ve|^2}$, and $\|\cdot\|_{*,\sigma,B_{\rho}}$ is defined in \eqref{starnorm}.
\end{prop}
\begin{proof}
We recast the system as a fixed point problem of the form\begin{equation}\label{fixedfinal}
\big(\mathbf{\Phi}_1,(d_{10},\dots,d_{N0})\big) =\widetilde{\mathcal{F}}\big(\mathbf{\Phi}_1,(d_{10},\dots,d_{N0})\big),
\end{equation}
where, for each $i=1,\dots,N,$ 
\begin{equation*}
\begin{aligned}
\phi_{i,1}&=T_i\left[H_i(\mathbf{\Phi}_1+\mathbf{\Phi}_2,\phi_{out})+\bar{B}_{i}[\phi_{i,2}]\right]\in \mathcal{X}_{*}\defeq\{\phi \in C^{2,\beta}(B_{\rho})\ : 
  \|\phi\|_{*, \sigma,B_{\rho}} <+\infty\},\\
d_{i0}&=\gamma_0\int_{\R^2} \left(H_i(\mathbf{\Phi}_1+\mathbf{\Phi}_2,\phi_{out})+\bar{B}_i[\phi_{i,2}]\right)Z_0, \quad \gamma_0^{-1}=\int_{\R^2}\gamma\Gamma^{\gamma-1}_{+}Z_0^2.
\end{aligned}
\end{equation*}
For clarity, we recall that $T_i$ is the linear operator of Proposition \ref{prop1}, \,$\phi_{i,2}(y)=d_{i0}\hat{\phi}_{2}(y)$ is given in \eqref{explicitphi2}, while Section \ref{subsecouter} provides the outer solution $\phi_{out}=\phi_{out}(\mathbf{\Phi}_1+\mathbf{\Phi}_2)$ through Proposition \ref{propouter}, satisfying \eqref{boundphiout} and \eqref{phioutP}. Moreover, the definitions of $H_i$ and $\bar{B}_i$ are introduced in \eqref{defofBinner} and \eqref{defHinner}.

By means of the Banach's fixed point Theorem, we seek a solution in the ball 
\begin{equation}\label{contractionball}
\begin{aligned}
\mathcal{B}=\Big\{\big(\mathbf{\Phi}_1,\left(d_{10},\dots,d_{N0}\right)\big)\in \mathcal{X}^{N}_{*}\times \R^{N}: &\|\phi_{i,1}\|_{*,\sigma,B_\rho}\leq \widehat{C}\ve\mu_i|\log\ve|,\\ &\,|d_{i0}|\leq \widehat{C}(\ve\mu_i)^{1+\sigma}|\log\ve|^{2+2\sigma},\,i=1,\dots,N\Big\},
\end{aligned}
\end{equation}
for some constant $\widehat{C}>0$ independent of $\ve$ and $\sigma^{*}$, to be chosen later.

 In this regard, we first need to show that $\widetilde{\mathcal{F}}\left(\mathcal{B}\right)\subset \mathcal{B}$. Assuming $\big(\mathbf{\Phi}_1,\left(d_{10},\dots,d_{N0}\right)\big)\in \mathcal{B},$ we obtain
\begin{equation*}
\begin{aligned}
\left|H_i\left(\mathbf{\Phi}_1+\mathbf{\Phi}_2,\phi_{out}\right)\right| &\leq C\Gamma^{\gamma-2}_{+}\bigg(\sum_{j=1}^{N}\left(|\tilde{\eta}_j\phi_{j,1}|^2+|\tilde{\eta}_j\phi_{j,2}|^2\right)+|\phi_{out}|^2\bigg) + |\widetilde{E}_i| +\left(\gamma\Gamma^{\gamma-1}_{+}+b_i(y)\right)|\phi_{out}|\\
&\leq C\ve\mu_i|\log\ve| \left(\Gamma\left(\frac{y}{2}\right)\right)^{\gamma-1}_{+},
\end{aligned}
\end{equation*}
thus we directly get \[\left\|H_i\left(\mathbf{\Phi}_1+\mathbf{\Phi}_2,\phi_{out}\right)\right\|_{2+\sigma,\beta,B_{\rho}} \leq C \ve\mu_i|\log\ve|.
\]
In addition, using \eqref{estimateforphi2} we find
\begin{equation*}
\begin{aligned}
\left|\bar{B}_i[\phi_{i,2}]\right| &\leq C\ve\mu_i\left(|y||D^{2}_{y}\phi_{i,2}|+|D_y \phi_{i,2}|\right)+\mathcal{O}\left(\ve\mu_i |\log\ve|\Gamma^{\gamma-2}_{+}\right)\phi_{i,2}\\
&\leq \frac{C\ve\mu_i|d_{i0}|}{2+|y|}+\mathcal{O}\left(\ve\mu_i|\log\ve|\Gamma^{\gamma-2}_{+}\right)|d_{i0}|\log(2+|y|),\\
\end{aligned}
\end{equation*}
so we further have
\[\left\|\bar{B}_i[\phi_{i,2}]\right\|_{2+\sigma,\beta,B_{\rho}}\leq C\ve\mu_i|\log\ve|.\]
Due to Proposition \ref{prop1}, we then deduce 
\begin{equation}\label{firsterr71}
\begin{aligned}
\left\|\phi_{i,1}\right\|_{*,\sigma,B_{\rho}}=\big\|T_i\left[H_i(\mathbf{\Phi}_1+\mathbf{\Phi}_2,\phi_{out})+\bar{B}_i[\phi_{i,2}]\right]\big\|_{*,\sigma,B_{\rho}} \leq C_1\ve\mu_i|\log\ve|,
\end{aligned}
\end{equation}
for some $C_1>0$.

To proceed, we obtain the estimate 
\begin{equation}\label{estimateBphi2Zo}
\begin{aligned}
\int_{B_{\rho}}\bar{B}_i[\phi_{i,2}]Z_0 
&\leq C\ve\mu_i  |d_{i0}|\int_{B_{\rho}}\frac{Z_0}{2+|y|}+ \int_{B_{\rho}}\mathcal{O}\left(\ve\mu_i|\log\ve|\Gamma^{\gamma-2}_{+}\right)|d_{i0}|\log(2+|y|) Z_0\\
&=\mathcal{O}\left(\frac{|d_{i0}|}{|\log\ve|}\right),
\end{aligned}
\end{equation}
and 
\begin{equation}\label{estissue}
\begin{aligned}
\int_{B_{\rho}}H_i\left(\mathbf{\Phi}_1+\mathbf{\Phi}_2,\phi_{out}\right)Z_0 &\leq C\int_{B_{\rho}}\Gamma^{\gamma-2}_{+}\bigg(\sum_{j=1}^{N}\big(|\tilde{\eta}_j\phi_{j,1}|^2+|\tilde{\eta}_j\phi_{j,2}|^2\big)+|\phi_{out}|^2\bigg)Z_0 +C\int_{B_{\rho}} \widetilde{E}_iZ_0\\
&\hspace{4mm}+C\int_{B_{\rho}}\left(\gamma\Gamma^{\gamma-1}_{+}+\mathcal{O}\big(\ve\mu_i|\log\ve|\Gamma^{\gamma-2}_{+}\big)\right)\phi_{out}Z_0.
\end{aligned}
\end{equation}
At this stage, we recognise that
\begin{equation}\label{estquad}
\begin{aligned}
\int_{B_{\rho}}\Gamma^{\gamma-2}_{+}\bigg(\sum_{j=1}^{N}\left(|\tilde{\eta}_j\phi_{j,1}|^2+|\tilde{\eta}_j\phi_{j,2}|^2\right)+|\phi_{out}|^2\bigg)Z_0&=\mathcal{O}\bigg(\sum_{j=1}^{N}\|\phi_{j,1}\|^2_{*,\sigma,B_{\rho}}\bigg)\\
&\leq C \ve^2\mu_i^2|\log\ve|^2.
\end{aligned}
\end{equation}
Nevertheless, even though $Z_0$ is a radial function, recalling the form of $\widetilde{E}_i$ from Proposition \ref{errorprop1}, namely 
\begin{align*}
\widetilde{E}_i&= \kappa_i\left(\frac{3R_ih^2+R_i^3}{2h\left(h^2+R_i^2\right)^{\frac3 2}}\right)\ve\mu_i y_1\Gamma^{\gamma}_{+} +\kappa_i\gamma\Gamma^{\gamma-1}_{+}\Bigg[\nu'(1)\log\left(1+\frac{\mu^{*}_i}{\mu^{0}_i}\right)+\left(c_{1,i}\ve\mu_iy_1+c_{2,i}\ve^2\mu_i^2|y|^2\right)\Gamma(y)\\
   &+\ve\mu_i y_1|\log\ve|\widehat{\mathcal{A}}_{1,i}(P)+\ve\mu_i y_2|\log\ve|\widehat{\mathcal{A}}_{2,i}(P)+\mathcal{O}\left(\ve^2\mu_i^2|\log\ve|^2 |y|^2\right)\Bigg]+\mathcal{O}\left(\ve^2\mu_i^2|\log\ve|^2|y|^2\Gamma^{\gamma-2}_{+}\right),
    \end{align*}
for the remaining integrals in \eqref{estissue} we obtain 

\begin{equation}\label{strong}
\int_{B{\rho}}\widetilde{E}_iZ_0 =\mathcal{O}\left(\ve^{1+\tilde{\sigma}}\right), \quad \int_{B_{\rho}}\gamma\Gamma^{\gamma-1}_{+}\phi_{out}Z_0=\mathcal{O}\left(\ve^{1+\sigma^{*}}\right),
\end{equation}
with $\tilde\sigma>0$ and $\sigma^{*}\in(0,\sigma)$ as in \eqref{scalingpert} and \eqref{boundphiout}, respectively.

As already alluded to in Remark \ref{decoupling}, the estimate for the second integral in \eqref{strong} is a manifestation of the strong coupling between the Inner and Outer problems, since the integrals in \eqref{estquad} and \eqref{strong} involving the inner perturbations $\mathbf{\Phi}_{1},\mathbf{\Phi}_2$ and the outer correction $\phi_{out}$ are of size \[\mathcal{O}\left(\ve^{1+\sigma^{*}}\right) \gg \mathcal{O}\left(\left(\ve\mu_i\right)^{1+\sigma}|\log\ve|^{2+2\sigma}\right)\]
for all small $\ve>0.$

To sufficiently decouple the system and ensure that $\widetilde{\mathcal{F}}(\mathcal{B})\subset\mathcal{B}$, we now make an explicit choice of the perturbations $\mu^{*}_i$ defining the scaling parameters $\mu_i=\mu_i^{0}+\mu_i^{*},\,i=1,\dots,N$ in \eqref{scalingansatz},  and verify a posteriori the estimate in \eqref{scalingpert}. We notice that in the bounded region
\[\mathcal{U}\defeq\Big\{x\in\R^2: |A_i^{-1}(x-P_i)|<\frac{2\delta_1}{|\log\ve|^2}\Big\},\] since $x-P_i=A_i\ve\mu_i y$, we can use the expansion
\begin{equation}\label{expphiout}
\phi_{out}(x)=\phi_{out}(P_i)+ \mathcal{O}\big(\ve\mu_i|A_iy|\|\nabla\phi_{out}\|_{L^{\infty}(\mathcal{U})}\big),
\end{equation}
where \eqref{boundphiout},\eqref{phioutP}, and a standard elliptic estimate give \[\|\nabla\phi_{out}\|_{L^{\infty}(\mathcal{U})}=\mathcal{O}(\ve^{1+\sigma^{*}}
),\] with $0<\sigma^{*}<\sigma$ as in \eqref{boundphiout}.

We then require
\begin{equation}\label{scalingpertexplicit}
-\kappa_i\nu'(1)\log\left(1+\frac{\mu_i^{*}}{\mu_i^{0}}\right)=\phi_{out}(P_i), \quad i=1,\dots,N,
\end{equation}
which implies $-\kappa_i\nu'(1)\log\left(1+\frac{\mu_i^{*}}{\mu_i^{0}}\right)=\mathcal{O}\left(\ve^{1+\sigma^{*}}\right)$ due to \eqref{boundphiout}, so that a posteriori $\tilde{\sigma}=\sigma^{*}$ in \eqref{scalingpert}. 

Making use of \eqref{expphiout} and \eqref{scalingpertexplicit} in \eqref{estissue} we get
\begin{equation}\label{finalestHiZ_0}
\int_{B_{\rho}}H_i\left(\mathbf{\Phi}_1+\mathbf{\Phi}_2,\phi_{out}\right)Z_0=\mathcal{O}\left(\ve^2\mu_i^2|\log\ve|^2\right),
\end{equation}
therefore combining \eqref{estimateBphi2Zo} and \eqref{finalestHiZ_0}, Proposition \ref{prop1} yields
\begin{equation}\label{contr2}
\begin{aligned}
|d_{i0}|&=\mathcal{O}\left(\frac{|d_{i0}|}{|\log\ve|}\right)+\mathcal{O}\left(\frac{\log(2+\rho)}{1+\rho^{\sigma}}\right)\left\|H_i\left(\mathbf{\Phi}_1+\mathbf{\Phi}_2,\phi_{out}\right)+\bar{B}_i[\phi_{i,2}]\right\|_{2+\sigma,\beta,B_{\rho}}\\
&\leq C_2 (\ve\mu_i)^{1+\sigma}|\log\ve|^{2+2\sigma},
\end{aligned}
\end{equation}
for some $C_2>0.$

Collecting \eqref{firsterr71} and \eqref{contr2}, we conclude that if $\widehat{C}>0$ in \eqref{contractionball} is chosen sufficiently large (independently of $\ve>0$ and $0<\sigma^{*}<\sigma$), then $\widetilde{\mathcal{F}}(\mathcal{B})\subset \mathcal{B}$. 

To complete the proof, it remains to establish that $\widetilde{\mathcal{F}}$ in \eqref{fixedfinal} is a contraction mapping in $\mathcal{B}$. In order to do so, we define 
\[\widetilde{\mathbf{\Phi}}^{\ell}=\sum_{i=1}^{N}\tilde{\eta}_i\left(\phi^{\ell}_{i,1}+\phi^{\ell}_{i,2}\right)+\phi^{\ell}_{out}, \quad\mbox{with} \quad \phi_{out}^{\ell}=\phi_{out}[\mathbf{\Phi}_{1}^{\ell}+\mathbf{\Phi}_{2}^{\ell}], \quad \ell=1,2,\]
and then set
\[H_i^{\ell}=H_i\left(\widetilde{\mathbf{\Phi}}^{\ell},\phi_{out}^{\ell}\right),\quad d^{\ell}_{i0}=\gamma_{0}\int_{R^2} \left(H_i^{\ell}+\bar{B}_i[\phi^{\ell}_{i,2}]\right)Z_0, \quad \gamma_0^{-1}=\int_{\R^2}\gamma\Gamma^{\gamma-1}_{+}Z_0^2.\]
Using the notation introduced above and assuming  $\big(\widetilde{\mathbf{\Phi}}^{\ell},\left(d_{10}^{\ell},\dots,d_{N0}^{\ell}\right)\big)\in\mathcal{B}$ for $\ell=1,2,$ we initially deduce from  \eqref{outeroperatorphiin} that
\begin{equation}\label{phioutcontraction}
\|\phi^1_{out}-\phi^2_{out}\|_{\infty}\leq C\bigg((\ve\mu)^{\sigma}|\log\ve|^{\zeta}\sum_{j=1}^{N}\|\phi^1_{j,1}-\phi^2_{j,1}\|_{*,\sigma,B_{\rho}}+|\log\ve|^{\zeta}\sum_{j=1}^{N}|d_{j0}^1-d_{j0}^2|\bigg),
\end{equation}
for some $\zeta>0$ and $\mu$ as in \eqref{maxscaling}.

Furthermore, since 
\begin{equation*}
      H_i^1-H_i^2=\mathcal{N}_i\left(\widetilde{\mathbf{\Phi}}^1\right)-\mathcal{N}_i\left(\widetilde{\mathbf{\Phi}}^2\right)+ \left(\gamma\Gamma^{\gamma-1}_{+}+b_i\right)\left(\phi_{out}^{1}-\phi_{out}^{2}\right),
\end{equation*}
we infer that
\begin{equation*}
\begin{aligned}
|H_i^1-H_i^2|&\leq C\Gamma^{\gamma-2}_{+}\bigg(\sum_{j=1}^{N}\tilde{\eta}_j\big(|\phi_{j,1}^1-\phi_{j,1}^2|^2+|\phi_{j,2}^1-\phi_{j,2}^2|^2\big)+|\phi^1_{out}-\phi^2_{out}|^2\bigg)\\
&\hspace{3mm}+\left(\gamma\Gamma^{\gamma-1}_{+}
+b_i\right)|\phi^1_{out}-\phi^2_{out}|,
\end{aligned}
\end{equation*}
where we obtain the estimate
\begin{equation}\label{Esth1minush2}
\begin{aligned}
\|H_i^1-H_i^2\|_{2+\sigma,\beta,B_{\rho}}&\leq C \Big(\|\phi_{out}^1-\phi_{out}^2\|_{\infty}+\sum_{j=1}^{N}\big(\|\phi^{1}_{j,1}-\phi^{2}_{j,1}\|^2_{*,\sigma,B_{\rho}}+|d^{1}_{j0}-d^{2}_{j0}|^2\big)+\|\phi_{out}^1-\phi_{out}^2\|^2_{\infty}\Big)\\
&\hspace{4mm}+C\ve\mu_i|\log\ve|\left(\sum_{j=1}^{N}\Big(\|\phi^{1}_{j,1}-\phi^{2}_{j,1}\|_{*,\sigma,B_{\rho}}+|d^{1}_{j0}-d^{2}_{j0}|\Big)+\|\phi_{out}^1-\phi_{out}^2\|_{\infty}\right).
\end{aligned}
\end{equation}
Moreover, due to \eqref{explicitphi2} and \eqref{estimateforphi2} we have \[\bar{B}_i[\phi^{1}_{i,2}-\phi^{2}_{i,2}]=\mathcal{O}(|d^{1}_{i0}-d^{2}_{i0}|)\bar{B}_i[\hat{\phi}_2], \quad \mbox{where} \quad \bar{B}_i[\hat{\phi}_{2}]\leq \frac{C\ve\mu_i}{2+|y|}+\mathcal{O}\left(\ve\mu_i|\log\ve|\Gamma^{\gamma-2}_{+}\right)\log(2+|y|),\]
hence we get
\begin{equation}\label{estforBfinal}
\left\|\bar{B}_i[\phi^{1}_{i,2}-\phi^{2}_{i,2}]\right\|_{2+\sigma,\beta,B_{\rho}} \leq \frac{\mathcal{O}\left(|d_{i0}^1-d_{i0}^2|\right)}{(\ve\mu_i)^{\sigma}|\log\ve|^{2+2\sigma}}.
\end{equation}
In addition, Proposition \ref{prop1} gives
\begin{equation}\label{finalestimates1}
\begin{aligned}
    |d^{1}_{i0}-d^{2}_{i0}|&\leq \mathcal{O}\left(\frac{\log(2+\rho)}{1+\rho^{\sigma}}\right)\left\|H_i^1-H_i^2+\bar{B}_i[\phi^{1}_{i,2}-\phi^{2}_{i,2}]\right\|_{2+\sigma,\beta,B_{\rho}}+\mathcal{O}\left(\frac{|d_{i0}^1-d_{i0}^2|}{|\log\ve|}\right)\\
    &\leq C(\ve\mu_i)^{\sigma}|\log\ve|^{1+2\sigma}\Bigg[\|\phi_{out}^1-\phi_{out}^2\|_{\infty}+\sum_{j=1
}^{N}\left(\|\phi^{1}_{j,1}-\phi^{2}_{j,1}\|^2_{*,\sigma,B_{\rho}}+|d^{1}_{j0}-d^{2}_{j0}|^2\right)\\
&\hspace{4mm}+\|\phi_{out}^1-\phi_{out}^2\|^2_{\infty}
+\ve\mu_i|\log\ve|\Big(\sum_{j=1}^{N}\big(\|\phi^{1}_{j,1}-\phi^{2}_{j,1}\|_{*,\sigma,B_{\rho}}+|d^{1}_{j0}-d^{2}_{j0}|\big)+\|\phi_{out}^1-\phi_{out}^2\|_{\infty}\Big)\Bigg]\\
&\hspace{4mm}+\mathcal{O}\left(\frac{|d_{i0}^1-d_{i0}^2|}{|\log\ve|}\right),
\end{aligned}
\end{equation}
while using \eqref{Esth1minush2}, \eqref{estforBfinal}, \eqref{finalestimates1}  and Proposition \ref{prop1} once again, we find 
\begin{equation}\label{finalestimates2}
\begin{aligned}
&\|\phi_{i,1}^1-\phi_{i,1}^2\|_{*,\sigma,B_{\rho}}=\big\|T_i\big[H_i^1-H_i^2+\bar{B}_i[\phi_{i,2}^1-\phi_{i,2}^2]\big]\big\|_{*,\sigma,B_{\rho}}\\
&\leq C\Bigg[\|\phi_{out}^1-\phi_{out}^2\|^2_{\infty}+\sum_{j=1}^{N}\Big(\|\phi^{1}_{j,1}-\phi^{2}_{j,1}\|^{2}_{*,\sigma,B_{\rho}}+|d^{1}_{j0}-d^{2}_{j0}|^2\Big)+\|\phi_{out}^1-\phi^2_{out}\|_{\infty}\\&\hspace{10mm}+\ve\mu_i|\log\ve|\Big(\sum_{j=1}^{N}\big(\|\phi^{1}_{j,1}-\phi^{2}_{j,1}\|_{*,\sigma,B_{\rho}}+|d^{1}_{j0}-d^{2}_{j0}|\big)+\|\phi_{out}^1-\phi^2_{out}\|_{\infty}\Big)\\
&\hspace{10mm}+\frac{1}{|\log\ve|}\Big(\sum_{j=1}^{N}\big(\|\phi^{1}_{j,1}-\phi^{2}_{j,1}\|^{2}_{*,\sigma,B_{\rho}}+|d^{1}_{j0}-d^{2}_{j0}|^2\big)+\|\phi_{out}^1-\phi_{out}^2\|_{\infty}\Big)\Bigg].
\end{aligned}
\end{equation}
Collecting together \eqref{phioutcontraction}, \eqref{finalestimates1} and \eqref{finalestimates2}, we  conclude that $\widetilde{\mathcal{F}}$ is a contraction mapping in $\mathcal{B}$ for all sufficiently small $\ve>0$, hence \eqref{fixedfinal} possesses a fixed point.
\end{proof}

\section{Balancing conditions and adjustment of the vortex locations}\label{reduced}
 In Section \ref{Section projected}, we have constructed a solution $(\phi_{in},\phi_{out})=(\phi_{in,1},\dots,\phi_{in,N},\phi_{out})$ to the coupled system 
\begin{equation*}
\Delta_y\phi_{in,i} +\gamma\Gamma^{\gamma-1}_{+}\phi_{in,i}+\bar{B}_i[\phi_{in,i}]+H_i(\phi_{in},\phi_{out})=\sum_{j=1}^{2}d_{ij}\gamma\Gamma^{\gamma-1}_{+}Z_j \, \inn \,\, B_{\rho}, \quad i=1,\dots,N,
\end{equation*}
and 
\[L_x[\phi_{out}]+\sum_{i=1}^{N} \left(L_x[\tilde{\eta}_i\phi_{in,i}]-\tilde{\eta}_iL_x[\phi_{in,i}]\right)=0 \, \inn \,\, \R^2.\]
For the reader's convenience, we recall that $\rho = \frac{2\delta_1}{\ve\mu_i|\log\ve|^2}$, $\bar{B}_i$ and $H_i(\phi_{in},\phi_{out})$ are defined in \eqref{defbarBinner}-\eqref{defHinner}, while the kernel elements $Z_j,\,j=1,2$ are given in \eqref{elemkernel}. Moreover, the solution  $(\phi_{in,1},\dots,\phi_{in,N},\phi_{out})$\,is obtained in Proposition \ref{propsec8} using a fixed point argument, and the corresponding estimates can be found in \eqref{boundsforsoln}.

To conclude the proof of Theorem \ref{maintheorem}, it remains to obtain an exact solution to the coupled Inner-Outer system in \eqref{inprobexp} and \eqref{poisreduced}, by solving   
\begin{equation}\label{zerofunctionals}
 d_{ij}\left[H_i(\phi_{in},\phi_{out})+\bar{B}_i[\phi_{i,2}]\right]=0, \quad \mbox{for all} \quad i=1,\dots,N, \, j=1,2,\end{equation}
where 
\begin{equation*}
        d_{ij}\left[H_i(\phi_{in},\phi_{out})+\bar{B}_i[\phi_{i,2}]\right]=\gamma_j \int_{\R^2}\left(H_i(\phi_{in},\phi_{out})+\bar B_i[\phi_{i,2}]\right)Z_j, \quad \gamma_{j}^{-1}=\int_{\R^2}\gamma\Gamma^{\gamma-1}_{+}Z_j^2.
        \end{equation*}
As will be demonstrated in the remainder of this section, satisfying these conditions requires a suitable choice of the points $P_1,\dots,P_N$ and $\tilde{r}$ in \eqref{formpoints}, thereby identifying the correct vortex locations for an exact solution of \eqref{euler}.

To begin with, Proposition \ref{prop1} together with \eqref{maxscaling} shows that $\|\bar{B}_i[\phi_{i,2}]+H_i(\phi_{in},\phi_{out})\|_{2+\sigma,\beta,B_{\rho}} \lesssim \ve\mu|\log\ve|,$ hence
\begin{equation*}
d_{ij}\left[H_i(\phi_{in},\phi_{out})+\bar{B}_i[\phi_{i,2}]\right]
=\gamma_{j}\int_{B_{\rho}} \left(H_i(\phi_{in},\phi_{out})+\bar{B}_i[\phi_{i,2}]\right)Z_{j}\,\dd y +\ve^{1+\sigma}\mathcal{M}(P),
\end{equation*}
for $\sigma>0$ as in \eqref{boundsforsoln}. 

Here and throughout this section, $\mathcal{M}(P)$ denotes a smooth function of $P=(P_1,\dots,P_N)$ satisfying \eqref{formpoints}, which is uniformly bounded as $\ve\to 0$. Its expression may vary from line to line, and even within the same line.

Furthermore, since $H_i(\phi_{in},\phi_{out})=\mathcal{N}_i\Big(\sum\limits_{j=1}^{N}\tilde{\eta}_j\phi_{in,j}+\phi_{out}\Big)+\left(\gamma\Gamma^{\gamma-1}_{+}+b_i\right)\phi_{out}+\widetilde{E}_i,$ one can also write
\[\int_{B_{\rho}} \bigg[\mathcal{N}_i\Big(\sum_{\ell=1}^{N}\tilde{\eta}_{\ell}\phi_{in,\ell}+\phi_{out}\Big)+\left(\gamma\Gamma^{\gamma-1}_{+}+b_i\right)\phi_{out}+\bar{B}_i[\phi_{i,2}]\bigg]Z_j =\ve^{1+\sigma}\mathcal{M}(P), \quad j=1,2.\]
Consequently, the conditions in \eqref{zerofunctionals} amount to establishing
\begin{equation}\label{conderror}
    \int_{B_{\rho}} \widetilde{E}_i Z_j=\ve^{1+\sigma}\mathcal{M}(P), \quad i=1,\dots,N,\,\, j=1,2.
\end{equation}
In light of Proposition \ref{errorprop1} and the expansion of the approximation error $\widetilde{E}_i$, the integral against $Z_1$ can be expressed as
\begin{equation*}
\begin{aligned}
   \int_{B_{\rho}}\widetilde{E}_iZ_1 \,\dd y= \ve\mu_i\kappa_i\Bigg[&\left(\frac{3R_i h^2+R_i^3}{2h(h^2+R_i^2)^{\frac 3 2}}\right)\int_{B_{\rho}}y_1\Gamma^{\gamma}_{+}Z_1 \, \dd y+c_{1,i}\int_{B_{\rho}}y_1\gamma\Gamma^{\gamma-1}_{+}\Gamma Z_1 \, \dd y\\
   &+\big(\mathcal{F}_{1,i}\left(P\right)+\log|\log\ve|\mathcal{M}(P)\big)\int_{B_{\rho}}\gamma\Gamma^{\gamma-1}_{+}y_1Z_1\,\dd y+ \ve\mu_i |\log\ve|^2\mathcal{M}(P)\Bigg],
\end{aligned}
\end{equation*}
where \[\mathcal{F}_{1,i}(P)=|\log\ve|\left(-c_{1,i}\nu'(1)-\frac{\alpha R_i h}{\kappa_i\sqrt{h^2+R_i^2}}\right)+\sum_{j\neq i}\frac{\kappa_j}{\kappa_i}\nu'(1)\frac{[A_j^{-1}(P_i-P_j)]_1}{|A_j^{-1}(P_i-P_j)|^2}.\]
Furthermore, upon testing with $Z_2$ we obtain
\[\int_{B_{\rho}} \widetilde{E}_i Z_2 \, \dd y=\ve \mu_i \kappa_i \Bigg[\mathcal{F}_{2,i}(P)\int_{B_{\rho}}\gamma\Gamma^{\gamma-1}_{+}y_2Z_2\,\dd y +\ve\mu_i|\log\ve|^2\mathcal{M}(P)\bigg],\]
with \[\mathcal{F}_{2,i}(P)=\sum_{j \neq i} \frac{\kappa_j}{\kappa_i}\nu'(1)\frac{[A_j^{-1}(P_i-P_j)]_2}{|A_j^{-1}(P_i-P_j)|^2}.\]
Using the form of the points $P_1,\dots,P_N$ in \eqref{formpoints}, namely
\begin{equation*}
P_i=(r_0+\tilde{r},0)+\frac{1}{|\log\ve|}\widehat{P}_{i},\qquad P_i=(a_i,b_i),\qquad \widehat{P}_i=(\hat{a}_i,\hat{b}_i) , \qquad i=1,\dots,N,
\end{equation*}
and defining the rescaled points
\[\widetilde{P}_i=(\tilde{a}_i,\tilde{b}_i)=\left(\frac{\sqrt{h^2+r_0^2}}{h}\hat{a}_i,\hat{b}_i\right),\]
a direct substitution in \eqref{conderror} together with \eqref{constantsj} and the expansion in \eqref{expinverse} yields the reduced problem 
\begin{equation}\label{Reduced1}
\begin{aligned}
    \left(-\frac{\nu'(1)R_i h}{2(h^2+R_i^2)^{\frac 3 2}}-\frac{\alpha R_i h}{\kappa_i\sqrt{h^2+R_i^2}}\right)+\sum_{j\neq i}\nu'(1)\frac{\kappa_j}{\kappa_i}\frac{[(\widetilde{P}_i-\widetilde{P}_j)]_1}{|(\widetilde{P}_i-\widetilde{P}_j)|^2}&=\frac{\log|\log\ve|}{|\log\ve|}\mathcal{M}(P),\\
    \sum_{j\neq i}\nu'(1)\frac{\kappa_j}{\kappa_i}\frac{[(\widetilde{P}_i-\widetilde{P}_j)]_2}{|(\widetilde{P}_i-\widetilde{P}_j)|^2}&=\frac{1}{|\log\ve|}\mathcal{M}(P).
\end{aligned}
\end{equation}
%where $\mathcal{M}(P)$ denotes once again a smooth function that is uniformly bounded as $\ve \to 0,$ for points $P=(P_1,\dots,P_N)$ satisfying \eqref{formpoints}.
We now reformulate the reduced system \eqref{Reduced1} in complex coordinates by identifying  $\widetilde{P}_j=\left(\widetilde{P}_{j,1},\widetilde{P}_{j,2}\right)$ with $\widetilde{P}_j=\widetilde{P}_{j,1}+i\widetilde{P}_{j,2},$ which allows us to recast it as a perturbation of the limit problem \eqref{balancingcond}, which reads 
\begin{equation}\label{limitproblem}
\sum_{j\neq i}\frac{\kappa_j}{\widetilde{P}_i-\widetilde{P}_j}=\left(\frac{\kappa_i h r_0}{2(h^2+r_0^2)^{\frac 3 2}}+\frac{\alpha h r_0}{\nu'(1)\sqrt{h^2+r_0^2}}\right), \quad i=1,\dots,N.
\end{equation}
With this in mind, our objective is to find a solution to \eqref{Reduced1} as a small perturbation of the solution of \eqref{limitproblem} by linearisation.

For $i=1,\dots,N,$ we define
\begin{equation}\label{deflhsrhs}
\mathbf{F}_i:\C^{N}\to \C, \quad \mathbf{F}_i(\widetilde{P})=\sum_{j \neq i}\frac{\kappa_j}{\widetilde{P}_i-\widetilde{P}_j},\quad\mbox{and}\quad \quad U_i(r)=\left(\frac{\kappa_i h r}{2(h^2+r^2)^{\frac 3 2}}+\frac{\alpha h r}{\nu'(1)\sqrt{h^2+r^2}}\right).
\end{equation}
Due to the specific choice of $\widetilde{P}^{b}$ in \eqref{balancingcond}, it follows that
\[\mathbf{F}_i\big(\widetilde{P}^{b}\big)=U_i(r_0),\]
while the derivative of $\mathbf{F}=\left(\mathbf{F}_1,\dots,\mathbf{F}_N\right)$ at $\widetilde{P}^{b}$ is given by
\begin{equation}\label{diffF}
\dd\mathbf{F}_{\widetilde{P}^{b}}=\begin{pmatrix}
-\sum\limits_{j\neq 1}\kappa_j T_{1j} & \kappa_2 T_{12} & \cdots & \kappa_N T_{1N} \\
\kappa_1 T_{21} & -\sum\limits_{j\neq 2}\kappa_j T_{2j} & \cdots & \kappa_N T_{2N} \\
\vdots & \vdots & \ddots & \vdots \\
\kappa_1 T_{N1} & \kappa_2 T_{N2} & \cdots & -\sum\limits_{j\neq N}\kappa_j T_{Nj}
\end{pmatrix},\quad T_{ij}=T_{ji}=\frac{1}{\left(\widetilde{P}_i^{b}-\widetilde{P}_j^{b}\right)^2}.
\end{equation}
At this point, we recognise that the vector $\hat{e}=(1,\dots,1)\in\C^{N}$ lies in $\operatorname{Ker}\left(\dd\mathbf{F}_{\widetilde{P}^{b}}\right)$, while the nondegeneracy condition on the point $\widetilde{P}^{b}$ is exactly the requirement that this kernel is one-dimensional, i.e. 
\[\operatorname{Ker}\left(\rm{d}\mathbf{F}_{\widetilde{P}^{b}}\right)=\mbox{span}\{\hat{e}\}.\]
In order to obtain a solution of \eqref{Reduced1} as a small perturbation of $\widetilde{P}^{b},$ we introduce ${\mathbf{q}}=({q}_1,\dots,{q}_N)\in\C^N$ and define $\tilde{\mathbf{q}}=(\tilde{q}_1,\dots,\tilde{q}_N)\in\C^{N}$ by 
\[\tilde{q}_j\defeq\operatorname{Re}\left(\frac{\sqrt{h^2+r_0^2}}{h}q_j\right)+i\operatorname{Im}(q_j),\quad j=1,\dots,N.\] Accordingly, we decompose
\[\widetilde{P}_j=\widetilde{P}^{b}_j+\tilde{q}_j=\operatorname{Re}\left(\frac{\sqrt{h^2+r_0^2}}{h}(P_j^{b}+q_j)\right)+i\operatorname{Im}\left(P_j^{b}+q_j\right), \quad j=1,\dots,N,\]
which in view of \eqref{deflhsrhs} allows us to rewrite the system in \eqref{Reduced1} as 
\begin{equation}\label{expreduced1}
\mathbf{F}_i(\widetilde{P})=U_i(R_i)+\hat{\sigma}_i, \quad i=1,\dots,N,
\end{equation}
with $\operatorname{Re}(\hat{\sigma}_i)=\frac{\log|\log\ve|}{|\log\ve|}\mathcal{M}(P)$ and $\operatorname{Im}(\hat{\sigma}_i)=\frac{\mathcal{M}(P)}{|\log\ve|}.$

We next consider the  expansions
\begin{align*}
\mathbf{F}(\widetilde{P})&=\mathbf{F}(\widetilde{P}^{b})+\dd\mathbf{F}_{\widetilde{P}^b}(\tilde{\mathbf{q}})+\mathcal{O}(|\tilde{\mathbf{q}}|^2),\\
U_i(R_i)&=U_i(r_0)+U^{'}_{i}(r_0)\left(\tilde{r}+\frac{h}{\sqrt{h^2+r_0^2}}\frac{\operatorname{Re}\left(\widetilde{P}_i^{b}+\tilde{q}_i\right)}{|\log\ve|}\right)+\mathcal{O}\left(\left(\frac{|\tilde{\mathbf{q}}|}{|\log\ve|}+|\tilde{r}|\right)^{2}\right)+\mathcal{O}\left(\frac{|\tilde{r}|}{|\log\ve|}\right),
\end{align*}
while calculating the derivative $U_i^{'}(r_0)$ and substituting the fixed value of $\alpha$ in \eqref{uniformspeed}, we deduce that \eqref{expreduced1} can be expressed as
\begin{equation}\label{exp2reduced}
\dd\mathbf{F}_{\widetilde{P}^{b}}(\tilde{\mathbf{q}})=\mathcal{J}\left(\tilde{r},\tilde{\mathbf{q}}\right)+\frac{\tilde{r}h}{2}\left(\frac{h^2-2r_0^2}{(h^2+r_0^2)^{\frac 5 2}}(\kappa_1,\dots,\kappa_N) -\frac{h^2}{\left(h^2+r_0^2\right)^{\frac 5 2}}\frac{\sum\limits_{i=1}^{N} \kappa_i^2}{\sum\limits_{i=1}^{N}\kappa_i}(1,\dots,1)\right),
\end{equation}
where $\mathcal{J}(\tilde{r},\tilde{\mathbf{q}})$ is a smooth function satisfying
\begin{equation}\label{propK}
    \mathcal{J}(0,0)=0, \quad \mathcal{J}_{\tilde{\mathbf{q}}}(0,0)=0, \quad D_{\tilde{r}}\mathcal{J}(0,0)=o(1)\to 0, \quad \mbox{as}\quad \ve\to 0.
\end{equation}
Moreover, in the admissible range of $\tilde{r}$ and $\tilde{\mathbf{q}}$ in \eqref{formpoints}, one finds $\operatorname{Re}\mathcal{J}(\tilde{r},\tilde{\mathbf{q}})=\mathcal{O}\left(\frac{\log|\log\ve|}{|\log\ve|}\right)$ and $\operatorname{Im}\mathcal{J}(\tilde{r},\tilde{\mathbf{q}})=\mathcal{O}\left(|\log\ve|^{-1}\right)$ as $\ve \to 0$. In addition, since $\dd\mathbf{F}_{\widetilde{P}^{b}}$ has a one-dimensional kernel, then the same is valid for $\left(\dd\mathbf{F}_{\widetilde{P}^{b}}\right)^{T}$. In fact, since $T_{ij}=T_{ji}$ in \eqref{diffF}, it is easy to verify that $\operatorname{Ker}\big((\dd \mathbf{F}_{\widetilde{P}^{b}})^{T}\big)=\mbox{span}\{\hat{\kappa}\}$, with \[\hat \kappa=(\kappa_1,\dots,\kappa_N)\in\C^N.\]
Now, we further note that the linear system in \eqref{exp2reduced} is solvable if the right-hand side is orthogonal to $\hat\kappa.$ Since its projection onto $\hat\kappa$ is equal to 
\[\left(\mathcal{J}(\tilde{r},\tilde{\mathbf{q}})\cdot\hat\kappa-
\tilde{r}h\frac{r_0^2}{(h^2+r_0^2)^{\frac 5 2}}\sum_{i=1}^{N}\kappa_i^2\right)\frac{\hat\kappa}{\sum\limits_{i=1}^{N}\kappa_i^2},\]
we initially study the projected problem
\begin{equation}\label{reducedproj}
\begin{aligned}
\dd\mathbf{F}_{\widetilde{P}^{b}}(\tilde{\mathbf{q}})&=\mathcal{J}\left(\tilde{r},\tilde{\mathbf{q}}\right)+\frac{\tilde{r}h}{2}\left(\frac{h^2-2r_0^2}{(h^2+r_0^2)^{\frac 5 2}}(\kappa_1,\dots,\kappa_N) -\frac{h^2}{\left(h^2+r_0^2\right)^{\frac 5 2}}\frac{\sum\limits_{i=1}^{N} \kappa_i^2}{\sum\limits_{i=1}^{N}\kappa_i}(1,\dots,1)\right)\\
&\hspace{3mm}-\left(\mathcal{J}(\tilde{r},\tilde{\mathbf{q}})\cdot\hat\kappa-
\tilde{r}h\frac{r_0^2}{(h^2+r_0^2)^{\frac 5 2}}\sum_{i=1}^{N}\kappa_i^2\right)\frac{\hat\kappa}{\sum\limits_{i=1}^{N}\kappa_i^2}.
\end{aligned}
\end{equation}
Since $\widetilde{P}^{b}$ is a nondegenerate solution, it follows from the properties in \eqref{propK} and the Implicit Function Theorem that \eqref{reducedproj} admits a unique local solution of the form $\tilde{\mathbf{q}}\defeq \tilde{\mathbf{q}}(\tilde{r})$ in a neighbourhood of $(\tilde{\mathbf{q}},\tilde{r})=(0,0)$. In particular, the estimates of $\mathcal{J}(\tilde r,\tilde{\mathbf{q}})$ imply that this solution satisfies  \[|\tilde{\mathbf{q}}(\tilde r)|\lesssim \frac{\log|\log\ve|}{|\log\ve|}.\]
Towards obtaining a genuine solution of \eqref{expreduced1}, the next step is to find a suitable  $\tilde{r}=\tilde{r}_{*}$ so that

\begin{equation}\label{fpbarr}
    \tilde{r}\mapsto \mathcal{J}\left(\tilde{r},\tilde{\mathbf{q}}(\tilde{r})\right)\cdot\hat\kappa-\tilde{r}h\frac{r_0^2}{(h^2+r_0^2)^{\frac 5 2}}\sum\limits_{i=1}^{N}\kappa_i^{2}=0.
\end{equation}
 Using once more the property
 $\operatorname{Re}\mathcal{J}(\tilde{r},\tilde{\mathbf{q}})=\mathcal{O}\left(\frac{\log|\log\ve|}{|\log\ve|}\right) $ as $\ve \to 0$, Banach's fixed point theorem provides a unique solution $\tilde{r}_{*}$ for \eqref{fpbarr}, satisfying \[|\tilde{r}_{*}|\lesssim \frac{\log|\log\ve|}{|\log\ve|}.\]
Collecting the aforementioned results, we deduce that the  solution of \eqref{expreduced1} is given by
\[
P_i=\left(
r_0+\tilde{r}_{*}+\frac{h}{\sqrt{h^{2}+r_0^2}}\frac{\operatorname{Re}\left(\widetilde{P}_i^{b}+\tilde{q}_i(\tilde{r}_{*})\right)}{|\log\ve|},
\;
\frac{\operatorname{Im}\big(\widetilde{P}_i^{b}+\tilde{q}_i(\tilde{r}_{*})\big)}{|\log\ve|}
\right)^{T},\quad i=1,\dots,N,
\]
which satifies the desired estimates in \eqref{pointstheorem}. This concludes the proof of Theorem \ref{maintheorem}.
\begin{section}{Appendix}
In this section, we provide the proofs of the linear inner theories stated in Section \ref{Section 7}. 
\begin{proof}[Proof of Lemma \ref{lemat}]
Let $y=re^{i\theta}, \, r=|y|$. We decompose $\phi$ and $h$ in Fourier series as 
$$ h(y) =  \sum_{k\in\Z}  h_k (r) e^{ik\theta} , \quad \phi(y)= \sum_{k\in\Z}\phi_{k} (r) e^{ik\theta},
$$
to obtain the ODEs  
\begin{equation}\label{01}
L_k [\phi_{k}] + h_k(r) = 0 , \quad r\in(0,\infty), \quad k \in \Z,
\end{equation}
with
$$
L_k\defeq \partial_{rr} + \frac 1r \partial_{r}- \frac{k^2}{r^2}+  \gamma\Gamma^{\gamma-1}_{+}  .
$$
For $k=0$, we have that $z_0(r) = \frac{2}{ \gamma-1}\Gamma(r)+r\Gamma'(r)$ satisfies $L_0[z_0]=0,$ hence if we denote by $\xi_0\in(0,1)$ the unique root of $z_{0}$, the function
$$
\phi_{0} (r) =     -z_0(r)\int_{\xi_0}^r  \frac {\rm d s}{ s z_0(s)^2}  \int_0^s h_0(\rho) z_0(\rho) \rho\, \dd \rho
$$
is a smooth solution of \eqref{01}. Observing that $\int_0^\infty h_0(\rho) z_0(\rho) \rho\, \dd\rho = \frac 1{2\pi}\int_{\R^2} h(y)Z_0(y)\, \dd y,$
we deduce  
$$
|\phi_{0}(r)| \, \le \,   C\Big[ \, \log (2 + r) \Big| \int_{\R^2} h(y)Z_0(y)\, \dd y      \Big|   \, +\,  (1+r)^{-\sigma} \|h\|_{2+ \sigma} \Big ].
$$
For $|k|=1$, one can verify that  $z_k(r)=-\Gamma'(r)$ satisfies $L_k[z_k]=0$, thus a smooth solution reads
$$
\phi_{k}  (r) =     \Gamma'(r)\int_r^\infty  \frac {\dd s}{ s\Gamma'(s)^2}  \int_0^s h_k(\rho) \Gamma'(\rho) \rho\, \dd\rho,
$$
with
$$
|\phi_{k}(r)| \, \le \,   C\Big[ \,  (1+ r) \sum_{i=1}^2 \Big| \int_{\R^2} h(y)Z_i(y)\, \dd y      \Big|   \, +\,  (1+r)^{-(1+\sigma)} \|h\|_{ 2+\sigma} \Big ].
$$
For $k=2$, there exists a function $z_2(r)$ satisfying $ L_2[z_2] = 0 $ and  $z_2(r)=\mathcal{O}(r^{2})$ as $r\to 0$ and $r\to \infty$, while for $|k|\ge 2$ we note that
$$
\widetilde\phi_{k}  (r) =   \frac 4{ k^2} z_2(r)\int_0^r  \frac {\dd s}{ s z_2(s)^2}  \int_0^s |h_k(\rho)| z_2(\rho) \rho\,\dd \rho
$$
is a positive supersolution for equation \eqref{01}. As a result, for $|k|\geq 2$ there exists a unique solution $\phi_{k}$  with $|\phi_{k}(r)| \le \widetilde\phi_{k}  (r)$, hence
$$
|\phi_{k}(r)| \, \le \,    \frac C {k^2}   (1+r)^{-\sigma} \|h\|_{2+\sigma}, \quad |k|\ge 2.
$$
Combining the above estimates for each mode, the function  $$\phi(y)= \sum_{k\in\Z} \phi_{k} (r) e^{ik\theta} $$  is a solution of \eqref{Linearisedinner} and defines a linear operator of $h$, satisfying 
\begin{equation}\label{cot} \begin{aligned}
 | \phi (y)|
\,  \le  \,  C \left [ \,  \log (2+|y|) \,\left|\int_{\R^2} h Z_0\right|  +    (1+|y|) \sum_{i=1}^2 \left|\int_{\R^2} h Z_i\right|  +  (1+|y|)^{-\sigma}\|h\|_{2+\sigma}   \,\right ]. \end{aligned}
\end{equation}
Under the further assumption  $\|h\|_{2+\sigma,\beta} <+\infty$, we next establish analogous estimates for first and second order derivatives of $\phi$. We let $R=|y|\gg 1$, fix the direction $\vec e=\frac{y}{|y|}\in\mathbb{S}^1$ so that $y=R\vec e$, and consider the rescaled functions
\[
\phi_{R}(\ell)  = {R^{\sigma}} \phi \left(y+R\ell\right), \quad h_R(\ell) = R^{2+\sigma}h \left(y+R\ell\right). \]
In the neighbourhood $U_{y}\defeq \{y+R\ell: |\ell|<\frac{1}{2}\}$, one finds
\[\Delta_\ell\phi_{R}(\ell)  +   R^2 \gamma\Gamma^{\gamma-1}_{+}(y+R\ell)\phi_{R}(\ell) + h_R (\ell)  = 0.\]
Setting \[m_i \defeq  \Big| \int_{\R^2}  h Z_i \Big|, \quad i=0,1,2,\]
we use \eqref{cot}, the fact that $\|h_{R}\|_{L^{\infty}\big(B_{\frac 1 2}(0)\big)}\leq C \|h\|_{2+\sigma},$ and a standard elliptic estimate to find
 \begin{equation}\label{holder1}
 \|\nabla_{\ell} \phi_{R} \|_{L^{\infty}\big( B_{\frac 14}(0) \big)} + \|\phi_{R} \|_{L^{\infty}\big( B_{\frac 12}(0) \big)} \, \le\,  C\Big [ m_0 R^{\sigma}\log R + \sum_{i=1}^2 m_i R^{1+\sigma} + \| h\|_{2+\sigma} \Big ].
 \end{equation}
In addition, since $[h_R(\ell) ]_{B_{\frac 12} (0),\beta}\leq C\| h\|_{2+\sigma,\beta}$, a combination of interior Schauder estimates and \eqref{holder1} gives
\begin{equation}\label{holder2}
\|D^2_{\ell} \phi_{R} \|_{L^\infty\big( B_{\frac 14}(0) \big) } +  [ D^2_{\ell} \phi_{R} ]_{B_{\frac 1 4}(0),\beta} \, \le\,  C\Big [ m_0 R^{\sigma} \log R +   \sum_{i=1}^2 m_i R^{1+\sigma} +   \| h\|_{2+\sigma,\beta} \Big ].
\end{equation}
Collecting together \eqref{cot}, \eqref{holder1} and \eqref{holder2}, the result of the lemma follows.
\end{proof}
\begin{proof}[Proof of Proposition \ref{Linearisedinner}]
As a preliminary step, we consider a standard linear extension operator $h\mapsto \tilde{h} $ to the whole $\R^2$,
so that the support of $\tilde{h}$ is contained in $B_{2\rho}$ and $\|\tilde{h}\|_{\sigma,\beta} \le C\|h\|_{\sigma,\beta, B_\rho}$, where $C>0$ is independent of $\rho$. In addition,
we assume that the coefficients of $\bar{B}_i$ belong in $C^1(\R^2)$, with compact support in $B_{2\rho}$.
 In this regard, we study the auxiliary projected problem in $\R^2$ given by

\begin{equation}\label{001} 
\Delta_y\phi_{i}  +  \gamma\Gamma^{\gamma-1}_{+}\phi_{i}  + \bar{B}_i[\phi_{i}]  + \tilde{h}(y)  = \sum_{j=0}^2  d_{ij}\gamma\Gamma^{\gamma-1}_{+} Z_j \inn \R^2,
\end{equation}
where $d_{ij} = d_{ij}[h,\phi_{i}]$ are the scalars defined as
\begin{equation}\label{coeffsforortho}
d_{ij}= \gamma_j \int_{\R^2} (\tilde{h}(y)+\bar{B}_i[\phi_{i}] )Z_j,  \quad \gamma_j^{-1} = \int_{\R^2}  \gamma\Gamma^{\gamma-1}_{+} Z_j^2, \quad j=0,1,2.
\end{equation}
For $j=1,2,$ since $Z_j=\mathcal{O}\left(\frac{1}{1+|y|}\right)$ as $|y| \to \infty,$ one finds 
\[\bar{B}_i[Z_j]\leq C\ve\mu_i\left(|y||D^2_{y}Z_j|+|D_y Z_j|\right) +\mathcal{O}\left(\ve\mu_i|\log\ve| \Gamma^{\gamma-2}_{+}\right)Z_j  \leq \frac{C\ve\mu_i|\log\ve|}{1+|y|^2},\]
thus we get
 $$
  \int_{\R^2} \bar{B}_i[\phi_{i}]Z_j =   \int_{\R^2} \phi_{i} \bar{B}_i[Z_j] \leq C\ve\mu_i|\log\ve|\|\phi_{i}\|_{*,\sigma},\quad j=1,2,
 $$
 where we used that $ \int_{\R^2} \frac {1}{1+|y|^{2+\sigma}}  <+\infty$.
 
On the other hand, since $Z_0(y)=\mathcal{O}\left(\log\left(2+|y|\right)\right)$ as $|y|\to\infty,$ for $j=0$ we instead have
 \[\bar{B}_i[Z_0]\leq \frac{C\ve\mu_i|\log\ve|}{1+|y|},\]
 hence
 \begin{equation*}
 \begin{aligned}
\int_{\R^2} \bar{B}_i[\phi_{i}]Z_0 &\leq C\ve\mu_i\int_{B_{2\rho}} \left(|y||D^2_{y}\phi_{i}|+|D_{y}\phi_{i}|\right)Z_0 +\int_{B_{2\rho}} \mathcal{O}\left(\ve\mu_i|\log\ve| \Gamma^{\gamma-2}_{+}\right)Z_0\\
&\leq C\left(\ve\mu_i\right)^{\sigma}|\log\ve|^{2\sigma-1}\|\phi_{i}\|_{*,\sigma}.
\end{aligned}
\end{equation*}
In addition, we observe that
\[\int_{\R^2\setminus B_{\rho}} \tilde{h}Z_0=\mathcal{O}\left(\frac{\log(2+\rho)}{1+\rho^{\sigma}}\right)\|h\|_{2+\sigma,\beta,B_{\rho}}, \quad \int_{\R^2\setminus B_{\rho}} \tilde{h}Z_j=\mathcal{O}\left(\frac{1}{1+\rho^{1+\sigma}}\right)\|h\|_{2+\sigma,\beta,B_{\rho}}, \, j=1,2,\]
and 
\[\bar{B}_i[\phi_{i}] \leq\frac{C\ve\mu_i|\log\ve|}{1+|y|^{1+\sigma}}\|\phi_{i}\|_{*,\sigma},\]
which further implies that  \[\big\| \bar{B}_i[\phi_{i}] \big\|_{2+\sigma,\beta} \leq \frac{C}{|\log\ve|}\|\phi_{i}\|_{*,\sigma,B_{\rho}}.\]
To solve \eqref{001}, we work in the Banach space $\mathcal{X}\defeq\{\phi_{i} \in C^{2,\beta}(\R^2)\ : 
  \|\phi_{i}\|_{*, \sigma} <+\infty\}$ and use the linear operator $\mathcal T$ of Lemma \ref{lemat} to reformulate the equation as 
\be\label{fp}
 \phi_{i}  =   \mathcal A  [\phi_{i}]  +  \mathcal {H},\quad \phi_{i} \in \mathcal{X},
\ee
where
$$
\mathcal A  [\phi_{i}]
 =  \mathcal T\Big [\bar{B}_i[\phi_{i}] -\sum_{j=0}^2 d_{ij}[0, \phi_{i}]\gamma\Gamma^{\gamma-1}_{+}Z_j  \Big] ,\quad
 \mathcal H  = \mathcal T\Big [  \tilde h   - \sum_{j=0}^2 d_{ij}[\tilde h,0] \gamma\Gamma^{\gamma-1}_{+}Z_j  \Big].
$$
Employing \eqref{coeffsforortho}, we obtain
\[\big\|\mathcal{A}[\phi_{i}]\big\|_{*,\sigma}\leq C\frac{\delta_1}{|\log\ve|}\|\phi_{i}\|_{*,\sigma},\quad \|\mathcal{H}\|_{*,\sigma}\leq C\|h\|_{2+\sigma,\beta,B_{\rho}},\]
hence the Contraction Mapping Theorem in $\mathcal{X}$ yields a unique solution to fixed point problem \eqref{fp} for all sufficiently small $\ve>0$, which defines a linear operator of $h$ and satisfies
$$
\|\phi_{i} \|_{*, \sigma} \ \le\  C \|h \|_{ 2+\sigma,\beta, B_{\rho}}.
$$
To this end, we set $T_i[h] = \phi_{i}\big|_{B_\rho}$ to complete the proof.
\end{proof}
\end{section}
\bibliographystyle{plain}
\bibliography{references}
\end{document}


